%%%%%%%%%%%%%%%%%%%%%%%%%%%%%%%%%%%%%%%%%%%%%%%%%%%%%%%%%%%%%%%%%%%%%%%%%%
%
% StMcEvent - User Guide and Reference Manual -- LaTeX Source
%
% $Id: StMcEvent.tex,v 2.10 2006/08/15 21:41:59 jeromel Exp $
%
% Authors:Michael A. Lisa
%         Thomas S. Ullrich
%         Manuel Calderon de la Barca Sanchez
%
%%%%%%%%%%%%%%%%%%%%%%%%%%%%%%%%%%%%%%%%%%%%%%%%%%%%%%%%%%%%%%%%%%%%%%%%%%
%
% Notes to the authors:
%
% - A template for a class reference is at the end of this file.
% - Wrap all names functions with \name{}
% - All code, examples, prototypes in \verb+ ... +\\
%   or \begin{verbatim} ... \end{verbatim}
% - Use \StMcEvent if you refer to the package itself (not the class)
%
% This file is best edit with xemacs and the 'Function' package loaded.
%
%%%%%%%%%%%%%%%%%%%%%%%%%%%%%%%%%%%%%%%%%%%%%%%%%%%%%%%%%%%%%%%%%%%%%%%%%%
%
% $Log: StMcEvent.tex,v $
% Revision 2.10  2006/08/15 21:41:59  jeromel
% Fix rhic -> rhic.bnl.gov
%
% Revision 2.9  2003/12/02 16:42:10  calderon
% fix a verb statement extending over two lines
%
% Revision 2.8  2003/10/08 20:27:17  calderon
% Update documentation on FTPC, CTB, BEMC hits.
%
% Revision 2.7  2003/05/15 18:29:04  calderon
% Added data members from modified g2t_event table:
% Event Generator Final State Tracks, N Binary Collisions,
% N Wounded Nucleons East and West, N Jets.
%
% Revision 2.6  2003/04/16 18:12:35  calderon
% Documentation on subprocess id.
%
% Revision 2.5  2000/06/06 02:59:11  calderon
% Introduction of Calorimeter classes.  Update documentation.
%
% Revision 2.4  2000/04/18 15:49:38  calderon
% Updated documentation on particle table, added Table of Contents for
% PDF.
%
% Revision 2.3  2000/04/18 00:55:47  calderon
% Updated documentation in class reference.  Added section about
% inclusion of particle table.
%
% Revision 2.2  2000/02/03 23:04:30  calderon
% Fixed typo that caused index to be generated in \texttt
%
% Revision 2.1  2000/02/03 03:36:59  calderon
% Documentation for version 2.0
%
% Revision 2.0  1999/11/17 02:01:07  calderon
% Completely revised for new StEvent
%
% Revision 1.4  1999/07/23 00:03:48  calderon
% Corrected documentation about macro location, and current libraries where StMcEvent runs
%
%
%%%%%%%%%%%%%%%%%%%%%%%%%%%%%%%%%%%%%%%%%%%%%%%%%%%%%%%%%%%%%%%%%%%%%%%%%%
\documentclass[twoside]{article}

\parindent 0pt
\parskip 6pt
\advance\textwidth by 80pt%
\advance\evensidemargin by -80pt%

\usepackage{graphicx}
\usepackage{psboxit}
\usepackage{amsmath}
\usepackage{amssymb}
\usepackage{amsfonts}
\usepackage{fancyhdr}
\usepackage{times}
\usepackage{verbatim}
\usepackage{makeidx}
\usepackage[dvips=true,hyperindex=true,colorlinks=true,linkcolor=blue,bookmarks=true]{hyperref}

\PScommands      % init boxit
\makeindex

%%%%%%%%%%%%%%%%%%%%%%%%%%%%%%%%%%%%%%%%%%%%%%%%%%%%%%%%%%%%%%%%%%%%
%
% Define header and footer style
%
%%%%%%%%%%%%%%%%%%%%%%%%%%%%%%%%%%%%%%%%%%%%%%%%%%%%%%%%%%%%%%%%%%%%
\pagestyle{fancyplain}
\rhead[\fancyplain{}{\bfseries\leftmark}]
      {\fancyplain{}{\bfseries\rightmark}}
\lhead[\fancyplain{}{\bfseries\rightmark}]
      {\fancyplain{}{\bfseries\leftmark}}
\rfoot[{}]{\fancyplain{}{\bfseries\thepage}}
\lfoot[\fancyplain{}{\bfseries\thepage}]{}
\cfoot{}

%%%%%%%%%%%%%%%%%%%%%%%%%%%%%%%%%%%%%%%%%%%%%%%%%%%%%%%%%%%%%%%%%%%%
%
% Typographic Conventions
%
%%%%%%%%%%%%%%%%%%%%%%%%%%%%%%%%%%%%%%%%%%%%%%%%%%%%%%%%%%%%%%%%%%%%
\newcommand{\name}[1]{\textsl{#1}}%  class-, function-, package names
\newcommand{\StEvent}{\textsf{StEvent}}
\newcommand{\StMcEvent}{\textsf{StMcEvent}}
\newcommand{\StAssociationMaker}{\textsf{StAssociationMaker}}

%%%%%%%%%%%%%%%%%%%%%%%%%%%%%%%%%%%%%%%%%%%%%%%%%%%%%%%%%%%%%%%%%%%%
%
% Define multiline labels for class reference
%
%%%%%%%%%%%%%%%%%%%%%%%%%%%%%%%%%%%%%%%%%%%%%%%%%%%%%%%%%%%%%%%%%%%%
\newcommand{\entrylabel}[1]{\mbox{\textbf{{#1}}}\hfil}%
\newenvironment{entry}
{\begin{list}{}%
    {\renewcommand{\makelabel}{\entrylabel}%
     \setlength{\labelwidth}{90pt}%
     \setlength{\leftmargin}{\labelwidth}
     \advance\leftmargin by \labelsep%
    }%
}%
{\end{list}}

\newcommand{\Entrylabel}[1]%
{\raisebox{0pt}[1ex][0pt]{\makebox[\labelwidth][l]%
    {\parbox[t]{\labelwidth}{\hspace{0pt}\textbf{{#1}}}}}}
\newenvironment{Entry}%
{\renewcommand{\entrylabel}{\Entrylabel}\begin{entry}}%
  {\end{entry}}

\begin{document}

%%%%%%%%%%%%%%%%%%%%%%%%%%%%%%%%%%%%%%%%%%%%%%%%%%%%%%%%%%%%%%%%%%%%
%
%    Title page
%
%%%%%%%%%%%%%%%%%%%%%%%%%%%%%%%%%%%%%%%%%%%%%%%%%%%%%%%%%%%%%%%%%%%%
\begin{titlepage}
\pagestyle{empty}
\vspace*{-35mm}
\begin{center}
  \mbox{\includegraphics[width=2cm]{StarIcon.eps}}
  {\Large\bf STAR Offline Library Long Writeup}
  \hfill\mbox{}\\[3cm]
  \mbox{\includegraphics[width=\textwidth]{StMcEventTitle.eps}}
  \hfill\mbox{}\\[3cm]
  {\LARGE User Guide and Reference Manual}\\[2cm]
  {\LARGE $  $}  \\[5mm] % replaced by cvs with current revision
  {\LARGE $  $}  % replaced by cvs with current revision
  \vfill
\end{center}
\cleardoublepage
\end{titlepage}
\pagenumbering{roman}

%%%%%%%%%%%%%%%%%%%%%%%%%%%%%%%%%%%%%%%%%%%%%%%%%%%%%%%%%%%%%%%%%%%%
%
%    Table of contents
%
%%%%%%%%%%%%%%%%%%%%%%%%%%%%%%%%%%%%%%%%%%%%%%%%%%%%%%%%%%%%%%%%%%%%
\tableofcontents
\cleardoublepage

%%%%%%%%%%%%%%%%%%%%%%%%%%%%%%%%%%%%%%%%%%%%%%%%%%%%%%%%%%%%%%%%%%%%
%
%    User Guide
%
%%%%%%%%%%%%%%%%%%%%%%%%%%%%%%%%%%%%%%%%%%%%%%%%%%%%%%%%%%%%%%%%%%%%
\pagenumbering{arabic}
\part{User Guide}
\clearpage

%%%%%%%%%%%%%%%%%%%%%%%%%%%%%%%%%%%%%%%%%%%%%%%%%%%%%%%%%%%%%%%%%%%%

\section{Introduction}

\StMcEvent\footnote{We will adopt the convention used in \StEvent\ 
    and write \StMcEvent\ when we refer to the package and
    \name{StMcEvent} when we refer to the class.} is a package that
compliments the use of \StEvent\ . The aim is for the user to be able
to access and analyze Monte Carlo with the same object oriented
approach as \StEvent\ . The top class is \name{StMcEvent}, and a
pointer to this  class enables access to all the relevant Monte Carlo
information.  Again, the pointer is obtained through a Maker, in this
case \name{StMcEventMaker}.  An example invocation is found in sec.~\ref{sec:howto}.

The information contained in \StMcEvent\ is designed to be simply
a reflection of the information already available in the existing
g2t tables.  The primary keys and foreign keys are replaced by
associations between classes through pointers, like in \StEvent\ .
At the moment, not all the information found on the g2t tables is
encapsulated in \StMcEvent\ , for a full description of what is
avaliable, please refer to sec.~\ref{sec:refman}.  Also note that
\StMcEvent\ is a work in progress, and if additional information
is needed, it can be added.

The goal of \StMcEvent\ is to be used in conjunction with \StEvent\ in
order to have an OO model for both pure Monte Carlo data and DST data
that has passed through the whole reconstruction chain.  To be able
to relate the 2 packages, there is an additional package that is
run after both \StEvent\ and \StMcEvent\ are filled: \StAssociationMaker\ .
\index{StAssociationMaker}
The aim of \StAssociationMaker\ is to establish the relationships
between the reconstructed and Monte Carlo information, so that
users can easily check whether a particular reconstructed hit or
track is found, and if that is the case,
to directly obtain the associated Monte
Carlo hit or track so that an assessment of the quality of the
data can be made.

\clearpage

%%%%%%%%%%%%%%%%%%%%%%%%%%%%%%%%%%%%%%%%%%%%%%%%%%%%%%%%%%%%%%%%%%%%

\section{How to read this document}

This document is divided in two parts, a user guide and a
reference manual. In the first part we rather concentrate on basic
questions and provide guidance to get started.  Some of the information
given here overlaps that one found in the \StEvent\ documentation,
and we provide it again here for completeness (We think that it is better
to be redundant than to have to remember which information is in which
manual).  The reference section
provides information on all available classes, their member functions
and related operators. The class references contain one or more
examples which demonstrate some features of each class and how to use
them.

New users should \textbf{not} start with the Reference section. It is
meant as a lookup when specific information is needed. Beginners
should study the User Guide and should make themselves familiar only
with those classes they will encounter more frequently:
\name{StMcEvent}  (sec.~\ref{sec:StMcEvent}),
\name{StMcTrack}  (sec.~\ref{sec:StMcTrack}),
\name{StMcVertex} (sec.~\ref{sec:StMcVertex}),
\name{StMcTpcHit} (sec.~\ref{sec:StMcTpcHit}).
\name{StMcSvtHit} (sec.~\ref{sec:StMcSvtHit}).
and \name{StMcFtpcHit} (sec.~\ref{sec:StMcFtpcHit}).



Understanding the various
examples certainly is the best way to get started. However, the examples
are given for illustration purposes only, and are not complete programs.

%%%%%%%%%%%%%%%%%%%%%%%%%%%%%%%%%%%%%%%%%%%%%%%%%%%%%%%%%%%%%%%%%%%%

\section{Further documentation}
\label{sec:furtherdoc}

\StMcEvent\ makes use of various classes from the StarClassLibrary (SCL).
To obtain the SCL documentation, the easiest thing to do is to go to
the web page set up by Gene van Buren. 
Go to
\begin{verbatim}
http://www.star.bnl.gov/STARAFS/comp/root/special_docs.html
\end{verbatim}
Here you can easily obtain the documentation several packages
(this one for example).
\index{Gene's Documentation Page}
\index{SCL} \index{StarClassLibrary}

%%%%%%%%%%%%%%%%%%%%%%%%%%%%%%%%%%%%%%%%%%%%%%%%%%%%%%%%%%%%%%%%%%%%

\section{Getting \StMcEvent\ Sources}  \index{Getting StMcEvent sources}

To access the complete source code proceed as follows:

\StMcEvent\ is under {\bf CVS} control at BNL.  It can
be accessed via \name{afs}: \index{afs} \index{CVS} \index{CVSROOT}
\begin{enumerate}
  \item Obtain an \name{afs} token: \name{klog -cell rhic}.
  \item Make sure \name{\$CVSROOT} is set properly:\\ %%$
    (i.e.~\name{CVSROOT = /afs/rhic.bnl.gov/star/packages/repository})
  \item Check-out package into your current working directory:\\
    \name{cvs checkout StRoot/StMcEvent}
\end{enumerate}

%%%%%%%%%%%%%%%%%%%%%%%%%%%%%%%%%%%%%%%%%%%%%%%%%%%%%%%%%%%%%%%%%%%%

\section{Coding Standards}
\index{Coding Standards}

\StMcEvent\ tries to follow the STAR coding guidelines as described on the
STAR $\rightarrow$ Computing $\rightarrow$ Tutorials $\rightarrow$ C++
coding standards web page: \\
http://www.star.bnl.gov/STAR/html/comp\_l/train/standards.html.\\
Here we
summarize the most relevant ones concerning the programmable interface:
\begin{itemize}
\item all classes, enumerations and functions start with the prefix
    \textbf{St}
\item all member functions start with a lowercase letter
\item header files have the extension \textbf{.hh}, source files have
    the extension\textbf{ .cc}
\item the use of underscores in names is discouraged
\item classes, methods, and variables have self-explanatory English
    names and first letter capitalization to delineate words
\end{itemize}

\StMcEvent\ also follows the following rules set forth in \StEvent\ 
\begin{itemize}
\item methods (member functions) which return a certain object have
    the same name as the referring object, i.e., without a preceding
    \name{get} prefix \\ (as in \name{StMcEvent::primaryVertex()})
\item methods which assign a value to a data member (or members) carry
    the prefix \name{set} followed by the name of the member (as in \name{StMcEvent::set\-Primary\-Vertex()}).
\item integer variables which serve as counter or indices and never
    can take negative values are consistently declared as
    \name{unsigned}.
\item Objects which are returned by pointer are not guaranteed to
    exist in which case a NULL pointer is returned. It is the user's
    responsibility to check for the return value and make sure the
    pointer is valid (non-zero) before she/he dereferences it.  Objects
    which are guaranteed to exist are returned by reference or by value.
\end{itemize}
%%%%%%%%%%%%%%%%%%%%%%%%%%%%%%%%%%%%%%%%%%%%%%%%%%%%%%%%%%%%%%%%%%%%
\section{Conventions}
\subsection{Numbering Scheme}
\label{sec:numberscheme}
\index{Numbering}
The information here is a reiteration of that found in the \StEvent\
Manual.  Both packages use the same conventions, but in order to
stress the difference between the official STAR numbering scheme
for subdetector components, and
the C/C++ conventions for indexing the containers with which we represent those
subdetectors, we repeat the cautions here.  The C/C++ convention
is that the first element in an array has the index 0.  This applies
to vectors, collections or other containers that allow indexing.
The TPC sectors and padrows,
the SVT layers, ladders and wafers, and the FTPC planes and sectors
all start counting at 1.  So it is important to remember this when
using the member functions that return the sector or padrow of a
TPC hit, for example, to address the elements in a container.  If
you want to obtain the hit container in which a particular hit
\verb+h+ is stored, you would do the following:

\begin{verbatim}
evt->tpcHitCollection()->sector(h.sector()-1)->padrow(h.padrow()-1).hits();
\end{verbatim}
Note the subtraction of 1, in order to index the array properly.
The functions that have this convention are:

\begin{itemize}
	\item \texttt{StMcTpcHit::sector()}
	\item \texttt{StMcTpcHit::padrow()}
	\item \texttt{StMcFtpcHit::plane()}
	\item \texttt{StMcSvtHit::layer()}
	\item \texttt{StMcSvtHit::ladder()}
	\item \texttt{StMcSvtHit::wafer()}
	\item \texttt{StMcSvtHit::barrel()}
	\item \texttt{StMcRichHit::pad()}
	\item \texttt{StMcRichHit::row()}
\end{itemize}

As of October 2003, the volume Id of the FTPC hits also contains a sector.  However the hits
are still only contained by a collection based on the plane.

\subsection {References and Pointers}
Part of the goal of \StMcEvent\ is the use of pointers and references
as opposed to id's and foreign keys.  However, the handling of
pointers must be done with care.
To reiterate wise words of caution: If a function returns an object by
\textit{reference},
the object is guaranteed to exist.  If a function returns an object by {\it pointer},
then it is a good idea to always check if the pointer is \verb+NULL+. Not doing
this can be the cause of many a headache and sleepless nights.  This is one of
the most common ``monsters under the bed'' that can painfully bite, 
and crash code all over the place.

\subsection{Units}
\index{units} \index{system of units}
\label{sec:units}

All quantities in \StMcEvent\ are stored using the official STAR units:
cm, GeV and Tesla.  In order to maintain a coherent system of units it
is recommended to use the definitions in \name{SystemOfUnits.h} from
the StarClassLibrary. They allow to 'assign' a unit to a given
variable by multiplying it with a constant named accordingly
(centimeter, millimeter, kilometer, tesla, MeV, ...).  The value of
the constants is thus that the result after the multiplication follows
always the STAR system of units.

The following example illustrates their use:
{\footnotesize
\begin{verbatim}
double a = 10*centimeter;
double b = 4*millimeter;
double c = 1*inch;
double E1 = 130*MeV;
double E2 = .1234*GeV;

//
//   Print in STAR units
//
cout << "STAR units:" << endl;
cout << "a = " << a << " cm" << endl;
cout << "b = " << b << " cm" << endl;
cout << "c = " << c << " cm" << endl;
cout << "E1 = " << E1 << " GeV" << endl;
cout << "E2 = " << E2 << " GeV" << endl;

//
//   Print in personal units
//
cout << "\nMy units:" << endl;
cout << "a = " << a/millimeter << " mm" << endl;
cout << "b = " << b/micrometer << " um" << endl;
cout << "c = " << c/meter << " m" << endl;
cout << "E1 = " << E1/TeV << " TeV" << endl;
cout << "E2 = " << E2/keV << " keV" << endl;
\end{verbatim}
}%footnotesize
The resulting printout is:
{\footnotesize
\begin{verbatim}
STAR units:
a = 10 cm
b = 0.4 cm
c = 2.54 cm
E1 = 0.13 GeV
E2 = 0.1234 GeV

My units:
a = 100 mm
b = 4000 um
c = 0.0254 m
E1 = 0.00013 TeV
E2 = 123400 keV
\end{verbatim}
}%footnotesize

Further documentation can be found in the StarClassLibrary manual
(see sec.~\ref{sec:furtherdoc}).

\subsection{Containers and Iterators}
\label{sec:containers}
\index{containers}
\index{iterators}

The containers used throughout \StMcEvent\
\begin{itemize}
  \item are STL vectors (see below)\ref{subsec:vectors}.
  \item store objects by pointer.
\end{itemize}
Because the are {\tt vector}s, they allow random access as in:

\verb+pointer_to_object = container[index];+

The containers are guaranteed to provide the following member functions:
\name{size()}, \name{begin()}, \name{end()}.  All collections in
\StMcEvent\ have a referring iterator defined.
The containers and iterators are all declared in the file
{\tt StMcContainers.hh}, for convenience.  In addition,
all the classes in \StMcEvent\ are {\tt \#include}d in
the header {\tt StMcEventTypes.hh}.
\index{{\tt \#include} files}
\index{include files}
\index{StMcEventTypes}
The names given to the containers are meant to reflect
the objects contained, as well as whether a container is
a {\it structural} container or not.
\index{structural containers}
By a structural container we mean a container that owns the objects
it contains.  This means that when a structural container is deleted,
the objects it holds get deleted as well.  The naming convention
is then:

\begin{itemize}
  \item A {\bf vec}tor of {\bf p}oin{\bf t}e{\bf r}s
will carry the prefix {\bf \tt StPtrVec}
  \item A {\bf s}tructural {\bf vec}tor of {\bf p}oin{\bf t}e{\bf r}s
will carry the prefix {\bf \tt StSPtrVec}
\end{itemize}

The name is completed with the type of the objects they contain.  So
the structural container of pointers to objects of type \name{StMcTpcHit} is
\name{StSPtrVecMcTpcHit} where we drop the \name{St} of the object being
contained.

In reality, whether a container is structural or not
is of little importance in terms of its usage, the interface and methods of
both are the same.  The distinction is made to keep the objects organized,
and to make sure everthing gets deleted, and deleted in only one place.  But
this mainly goes on behind the scenes.  The main point one should
keep in mind from all of this is that, unless you really know
what you are doing, you should {\bf NOT} delete a
structural container.  (Play music for ``The Twilight Zone''.)


In order to keep your code independent of the underlying container
types, people are encouraged to use \textit{iterators} instead of
indices. These permit a higher degree of flexibility,
allowing to change containers
without changing the application code using them. The following
example demonstrates this:

Given an arbitrary \StMcEvent\ collection \name{anyColl} of type
\name{StAnyColl}
which holds objects of type \name{obj} the following code is
not guaranteed to work:
\begin{verbatim}
     for (int i=0; i<anyColl.size(); i++)
         obj = anyColl[i];
\end{verbatim}
A {\tt list} for example, does not allow random access, so indexing
is not supported for {\tt lists}.  For vectors and deques, indexing does work.
Iterators however, work for all STL containers.  The code above then,
should be replaced by:
\begin{verbatim}
     StAnyCollIterator iter;
     for (iter = anyColl.begin(); iter != anyColl.end(); iter++)
         obj = *iter;
\end{verbatim}

The names of the iterators defined in \StMcEvent\ are just the name
of the object contained + the word {\tt Iterator}.  For example,
the iterator for the container \name{StPtrVecMcVertex} is
\name{StMcVertexIterator}.

Note that all collections are by pointer, so they are
\textit{polymorphic} containers; that is, they are not restricted to
collect objects of the base type only, but can store
objects of any type derived from
the base. This is especially important when dereferencing the iterator to
access the objects in the collection.  Polymorphism is
not used in \StMcEvent (at least not as much as it could be used), but
it is necessary in \StEvent, where,
for example, one has containers of tracks, but they can be Global
or Primary.  If a method returns a base class and one really wants
a derived class, we have to use Run Time Type Information (RTTI),
{\tt dynamic\_cast}ing and such.  Also, make sure you know how to
get the objects you want, depending on the type of objects in the
container.

For a \textit{by-pointer} collection:
\begin{verbatim}
     StMcVertexIterator iter;
     StSPtrVecMcVertex& vertices = event->vertices();
     StMcVertex* vertex;
     for (iter = vertices.begin(); iter != vertices.end(); iter++)
         vertex = *iter;                         // dereference once
\end{verbatim}

Hint: Use pointers or references to collection elements wherever
possible.  Making local copies is often time- and memory-intensive.
For example the same code above but written as:
\begin{verbatim}
     StMcVertexIterator iter;
     StMcVertexCollection* vertices = event->vertexCollection();
     StMcVertex vertex;
     for (iter = vertices->begin(); iter != vertices->end(); iter++)
         vertex = **iter;                         // dereference twice
\end{verbatim}
invokes the \name{StMcVertex} assignment operator and creates a local copy.
This method is only useful if you intend to modify an object locally
but want to leave the original untouched.

\subsubsection{Quick glance at {\tt vector}s}
\label{subsec:vectors}
As mentioned before, all containers used in \StMcEvent\ are STL {\tt vectors}.  That
means that one can use the methods supported by the {\tt vector}
template class.  I'll give a brief overview of some basic elements
of {\tt vector}s to provide some common ground on which to walk on, for
people unfamiliar with the STL.  A LOT more information can be found
in your favorite C++ book.
\index{vector}

\paragraph{Common Member Types}
The type of the elements of the container is passed as the
first template argument, and is known as its {\tt value\_type}.  For
example, the {\tt value\_type} of the \name{StPtrVecMcTrack} container,
which is really a \verb+vector<StMcTrack*>+ is, of course, {\tt StMcTrack*}.
The type used for indexing into the container is known as {\tt size\_type},
and {\tt difference\_type} is the type of the result of subtracting
two iterators.  Like mentioned before, every container defines an
{\tt iterator} and a {\tt const\_iterator} for pointing to the elements
of the container.  One of the aims of these types is to allow the
possibility to write code without knowing the actual types involved.

\paragraph{Public Member Functions}
The following methods illustrate some of the methods provided by
the {\tt vector} container.  The list is by no means complete, but
even so there are some more methods listed here than normally used.
The counterpart {\tt const} methods are omitted.
\begin{Entry}
\item[Iterator Methods]
  \verb+iterator  begin()+\\
  Points to the first element of the container.

  \verb+iterator  end()+\\
  Points to the last-plus-one element.

  \verb+reverse_iterator  rbegin()+\\
  Points to the first element of the reverse sequence.
  Note that the
  type of the iterator is different than for {\tt begin()}.
  Useful when going through a sequence backwards.  For more
  information on {\tt reverse\_iterator} consult your favorite
  C++ reference book.
  
  \verb+reverse_iterator  rend()+\\
  Points to the last-plus-one element of the reverse sequence.

\end{Entry}

The following methods allow to access the elements of the container.

\begin{Entry}

\item[Useful Access \\Methods]

  \verb+reference  operator[](size_type n)+\\
  Provides unchecked access to the elements of the container.  Safe
  to use when there is a previous condition to guarantee that
  the argument is within bounds.

  \verb+reference  at(size_type n)+\\
  Provides range checked access to the elements of the container.\\
  Throws {\tt out\_of\_range} if the index is out of range.

  \verb+reference  front()+\\
   Reference to the first element of the container.
  
  \verb+reference  back()+\\
   Reference to the last element of the container.

  \verb+size_type  size()+\\
   Number of elements in the container.  

  \verb+void  push_back(const T& x)+\\
  Adds the element to the end of the container, {\tt size()} increases
  by 1.

  \verb+void  pop_back(const T& x)+\\
  Removes the last element of the container, {\tt size()} decreases
  by 1.

  \verb+void  clear()+\\
  Erases all the elements.

\end{Entry}


%%%%%%%%%%%%%%%%%%%%%%%%%%%%%%%%%%%%%%%%%%%%%%%%%%%%%%%%%%%%%%%%%%%%
\subsection{Particles from Event Generators}
\label{sec:particle}
\index{particle table}
\index{event generator tracks}

\StMcEvent also provides information on particles coming from the so
called ``particle table''.  This is where all the particles produced
by the event generator are stored, with all the necessary bookkeeping.

The particle table is useful for the cases where a particle that one is
interested in, has been decayed by the event generator and not in
GSTAR.  It will therefore not appear in the g2t\_track table.  In
order to find out if a track from the g2t\_track table comes from
our particle of interest in this case, we need to access the particle
table.

All entries in the particle table are kept in the same container
as the tracks from the g2t\_track table.  (We probably shouldn't call
them tracks because these particles will most definitely NOT leave a
track in the detector simulation, but we keep them along with the others
for convenience.)  There will of course be tracks that appear in both
tables, namely the final state particles of the event generator.  This
case is already accounted for in \StMcEvent, so that only one entry
with the relevant information from both tables is kept in the container.

To properly use the additional information, it is important to know how to
distinguish an \name{StMcTrack} that comes from the particle table
and one that doesn't.  This can be done in several ways, according to
the following properties (or conventions, if you will):

\begin{itemize}
	\item {\bf Tracks with an event generator label $> 0$ have a corresponding
	entry in the particle table.}
	\item {\bf Tracks with a key $> 0$ have a corresponding entry in the g2t\_track table.}
	\item {\bf Only tracks from the g2t\_track table are assigned a start vertex.} This also means that
	the method \name{StMcVertex::daughters()} will {\bf NOT} return any tracks from the particle
	table.  It is important to realize this if one wants to asks for the daughters of the
	primary vertex.  This will only return the tracks that actually have the primary
	vertex as their parent, according to the start\_vertex\_p entry of g2t\_track the table.
	(These tracks
	are most probably also in the particle table, this scheme will of course not exclude them.  But
	the point is that tracks that only appear in the particle table will {\bf NOT} have a start vertex
	and they will not appear as daughters of any vertex, {\it not even the primary}).
	\item {\bf To find out the genealogy of a particle, the \name{StMcTrack::parent()} method is provided.}
	This method will return a pointer to the parent track, if there is one, otherwise it will return
	a null pointer.  Using this method, one can get to the parent particle of a decay daughter all
	the way into the particle table.  The parentage is followed not just until we reach the
	particle table, but it is also followed {\it within} the particle table, i.e. down to the quarks and gluons
	of the event generator.
	It should be stressed, of course, that
	this further parentage is {\it only} event generator bookkeeping, and should not be used for more than this.
	For g2t\_tracks, one can also find daughters by going through their stop vertex (if there is one),
	which then knows about its daughter tracks.
\end{itemize}
Let's illustrate this with a simple example.  Suppose that we have an event generator and
we want to look at $\phi \longrightarrow e^{+}e^{-}$ and the $\phi$ was decayed by the
event generator.  To get to them, we can do the following.
\begin{enumerate}
	\item We get a pointer to the primary vertex (\name{StMcEvent::primaryVertex()}).
	\item We get the daughters of the primary vertex (\name{StMcVertex::daughters()}).
	\item We loop over the primary tracks thus obtained (Look in the containers section ~\ref{sec:containers}).
	\item For each of primary tracks, we check whether it is an electron (\name{StMcTrack::particleDefinition()}).
	\item If it is, we check whether its parent is a $\phi$ (\name{StMcTrack::parent()}, followed by
	\name{StMcTrack::particleDefinition()}.  Be careful with NULL pointers!).
	\item If it is a $\phi$, then we found it and we can do whatever we want with it, e.g. fill a 2-D
	histogram of $p_{\bot}$ vs. $y$.
\end{enumerate}
For more information on the specific methods to do this, refer to the relevant sections for
\name{StMcVertex} ~\ref{sec:StMcVertex} and \name{StMcTrack}  ~\ref{sec:StMcTrack}.

%%%%%%%%%%%%%%%%%%%%%%%%%%%%%%%%%%%%%%%%%%%%%%%%%%%%%%%%%%%%%%%%%%%%

\section{How to use \StMcEvent: The Macro and the Maker}
\label{sec:howto}
\subsection{StMcEventReadMacro.C}
\index{StMcEventReadMacro.C}
\index{root4star}
\index{ROOT files}
\textbf{[Note]} The example macro \texttt{StMcEventReadMacro.C} described in
this section reflects the original 1999--2006 CVS layout and ROOT~3/4 API.
Current usage examples can be found in
\texttt{StRoot/StMcEventMaker/macros/} and the STAR software documentation
wiki at \texttt{https://drupal.star.bnl.gov/STAR/comp}.

\StMcEvent\ has now been a package in the STAR library for a while.  The
new version which this guide refers to is, as of this writing,
available in starnew and stardev.  Remember, that libraries
get moved, so make sure you are using the right library for
your purposes.  This is very important, since to use \StMcEvent\
along with \StEvent\ one CANNOT combine the latest code with
files produced before December 1999, since the dst tables
and \StEvent\ changed substantially.

To run the example macro StMcEventReadMacro.C, one can just
run the default version directly from the library:
\begin{verbatim}
    stardev
    cd somewhere
    root4star -b
    .x StMcEventReadMacro.C
\end{verbatim}
This will run the macro in the repository with the default settings:
Process one event from the default geant.root file, load StMcEvent
and print out to the screen some information on the event.
If this step doesn't work, then something weird is going on.  (Although
make sure if you are using dev that Lidia is not rebuilding the whole
library while you are trying to run one of the examples!)  The
default macros are meant to work ``out of the box'', so to speak,
so this is the first step to check that things are working
properly.

The next step is to actually check out the macro and edit it yourself for
your purposes.  To check out the macro, do the following:
\begin{verbatim}
    klog
    mkdir workdir
    cd workdir
    cvs co StRoot/macros/examples/StMcEventReadMacro.C
\end{verbatim}

StMcEventReadMacro.C is found in
\name{\$CVSROOT/StRoot/macros/examples/StMcEventReadMacro.C},  %$
and it runs a chain runs a chain of just one maker to load \StMcEvent .


%%%%%%%%%%%%%%%%%%%

At the moment, it runs the chain on a ROOT file tree.  That is, it uses
the GEANT branches of the files it finds in the directory.  This is important to
realize if you want to run the macro on files other than the default, an
activity one should quickly graduate to.  The macro takes as arguments the
number of events to process, and the name of the file to be used.  An example invocation
is:

{\footnotesize

\begin{verbatim}

.x StMcEventReadMacro.C(10,"/star/rcf/test/dev/tfs_Linux/Mon/year_2a/hc_standard/*.geant.root")
     

\end{verbatim}
}
This invocation will process 10 events from the specified file.
Now, to open the files, branches, etc. we rely on StIOMaker.  This maker
takes the path to look for files from the specified file, but the files
it actually
opens depend on
what branches are activated in the macro.
This is done inside the macro with the command \name{SetBranch}.

\subsection{StMcEventMaker}
\index{StMcEventMaker}
\textbf{[Note]} \name{StMcEventMaker} is maintained in the separate package
\texttt{StRoot/StMcEventMaker}, not in this package.
The description below reflects the original design; consult the current
source in \texttt{StRoot/StMcEventMaker/} for the up-to-date interface.

\name{StMcEventMaker} is the code that actually does the opening of
the geant.root file and uses the information there to set up the
whole \StMcEvent .  If you want to use \StMcEvent\ you have to make
sure that \name{StMcEventMaker} is instantiated in the chain you
are running.  This will take care of all the setup.  The data
members and methods of the maker that you will find useful are:

\index{StMcEventMaker|textbf}
\label{sec:StMcEventMaker}
\begin{Entry}
\item[Synopsis]
    \verb+#include "StMcEventMaker.hh"+\\
    \verb+class StMcEventMaker;+\\

\item[Public Data\\ Members]
    \verb+Bool_t  doPrintEventInfo;+\\
    Print or do not print info on the current \StMcEvent\ event.
    (default=kFALSE).  This produces a lot of output, and is
    meant for debugging.  Every major
    class is dumped, the sizes of all collections, and the first
    element in every container. Don't use it for production.
    
    \verb+Bool_t  doPrintMemoryInfo;+\\
    Switch on/off checks on memory usage of \StMcEvent\
    (default=kFALSE).  In order to get a memory snapshot we use
    \texttt{StMemoryInfo} from the \name{StarClassLibrary}.  A
    snapshot is taken before and after the setup of \StEvent.  The
    numbers in brackets refer to the difference. Not available on SUN
    Solaris yet.
    
    \verb+Bool_t  doPrintCpuInfo;+\\
    Switch on/off CPU usage (default=kFALSE). Tells you how long it
    took to setup \StEvent. Timing is performed using \texttt{StTimer}
    from the StarClassLibrary.

    \verb+Bool_t  doUseTpc;+\\
	Load Tpc hits in the run (default = true).  If you do not want to
	look at Tpc hits, switch this off.

    \verb+Bool_t  doUseSvt;+\\
	Load Svt hits in the run (default = true).  If you do not want to
	look at Svt hits, switch this off.

    \verb+Bool_t  doUseFtpc;+\\
	Load Ftpc hits in the run (default = true).  If you do not want to
	look at Ftpc hits, switch this off.

    \verb+Bool_t  doUseRich;+\\
	Load Rich hits in the run (default = true).  If you do not want to
	look at Rich hits, switch this off.

    \verb+Bool_t  doUseBemc;+\\
	Load Bemc hits in the run (default = true).  If you do not want to
	look at Bemc hits, switch this off.  This loads the barrel towers.

    \verb+Bool_t  doUseBsmd;+\\
	Load Shower Max hits in the run (default = true).  This loads hits from
	Bsmde, and Bsmdp detectors.

    \verb+Bool_t  doUseCtb;+\\
	Load CTB hits in the run (default = true).  If you do not want to
	look at CTB hits, switch this off.

    \verb+Bool_t  doUseTofp;+\\
	Load TOFp hits in the run (default = true).  If you do not want to
	look at TOFp hits, switch this off.

    \verb+Bool_t  doUseTof;+\\
	Load TOFr hits in the run (default = true).  If you do not want to
	look at TOFr hits, switch this off.

    \verb+Bool_t  doUsePixel;+\\
	Load Pixel hits in the run (default = true).  If you do not want to
	look at Pixel hits, switch this off.  The pixel is under development
        (August 2003), so Pixel hits will not be found in general purpose
        simulation files for a while.

\item[Public Member\\ Functions]
    \verb+StMcEvent*  currentMcEvent();+\\
    Returns a pointer to the current \name{StMcEvent} object.
\end{Entry}

Since the \name{StMcEvent} object you get when you invoke
{\tt currentMcEvent()} is held by pointer, dont forget to
check if this is {\tt NULL}. This would immediately signal
something gone awry.  Also, do not delete the objects
you get from the StMcEvent methods, they will be deleted
every time you read a new event.

In order to use the classes in
\StMcEvent\ in other Makers, one header file with {\it all} the include files
is provided: {\tt StMcEventTypes.hh}. \index{StMcEventTypes}
Like with \StEvent\ this makes
life easier in terms of not having to remember which file you need where,
you just care about one include file and that's it.  However, convenience
comes at a price, because if you \#include all the files, you will certainly
be creating dependencies on classes you won't need or care about.  So if this
is an issue, then proceed as usual and only \#include the classes that you actually
do depend on.
\index{{\tt \#include} files}
\index{include files}
The next step is the use of the \StAssociationMaker\ package.  This is described
in its own manual.

%%%%%%%%%%%%%%%%%%%%%%%%%%%%%%%%%%%%%%%%%%%%%%%%%%%%%%%%%%%%%%%%%%%%

\section{Known Problems and Fixes}
\index{known problems and fixes}

\StMcEvent\ is changing as the software environment changes.  From the last
version, several things have been updated and the design revamped to
follow \StEvent\ .
The creation of the vertex collection from the tables took its final
form, once the g2t\_vertex table was corrected.  The g2t\_event table was
finally written out, so it is now available.

The FTPC and SVT associations between reconstructed and Monte Carlo
hits are now included.
Also, the problems \StAssociationMaker\ had with the declaration of multimaps
by the Solaris CC4.2 compiler (i.e. no Template Default Arguments, and
ObjectSpace implementation of the STL) have been solved and
now \StMcEvent\ and \StAssociationMaker\ compile and run
on Linux gcc, Solaris CC4.2 and CC5 and HP aCC.

The particle table has also been added to \StMcEvent, and this
allows to find particles of interest that have been decayed by
the event generator.

The switches to turn on/off the loading of the different detector hits
have been added for convenience.  The first version of the classes
for the EMC hits implemented by Aleksei Pavlinov are now also included.

The TOF classes have been added in Aug 2003, as well as the development
classes for the Pixel detector.

For corrections to this manual, typos, suggestions, etc. send an email
to

{\tt calderon@star.physics.yale.edu}.
\clearpage

%%%%%%%%%%%%%%%%%%%%%%%%%%%%%%%%%%%%%%%%%%%%%%%%%%%%%%%%%%%%%%%%%%%%
%
%    Reference Manual
%
%%%%%%%%%%%%%%%%%%%%%%%%%%%%%%%%%%%%%%%%%%%%%%%%%%%%%%%%%%%%%%%%%%%%
\part{Reference Manual}
\label{sec:refman}
\clearpage

%%%%%%%%%%%%%%%%%%%%%%%%%%%%%%%%%%%%%%%%%%%%%%%%%%%%%%%%%%%%%%%%%%%%

\section{Global Constants}

We use the constants defined in the two header files
\index{StarClassLibrary} \name{SystemOfUnits.h} and
\name{PhysicalConstants.h} which are part of the StarClassLibrary.
The types defined therein are used
throughout \StMcEvent\ .



%%%%%%%%%%%%%%%%%%%%%%%%%%%%%%%%%%%%%%%%%%%%%%%%%%%%%%%%%%%%%%%%%%%%

\section{Class Reference}
The classes which are currently implemented and available from the
STAR CVS repository are described in alphabetical order.

Inherited member functions and operators are not described in the
reference section of a derived class. Always check the section(s)
of the base class(es) to get a complete overview on the available
methods.

Note that some constructors are omitted, especially the ones which
take tables as arguments. They are for internal use only (even if
public).

Destructors, assignment operators and copy constructors are not
listed.  Macros and \name{Inline} declarations are omitted throughout the
documentation, as well as the {\tt virtual} keyword.  For the
most up-to-date reference of what is available, there is no
substitute to looking directly at the class definition in
the header file.

\clearpage

%%%%%%%%%%%%%%%%%%%%%%%%%%%%%%%%%%%%%%%%%%%%%%%%%%%%%%%%%%%%%%%%%%%%
%
%    Reference: StMcCalorimeterHit
%
%%%%%%%%%%%%%%%%%%%%%%%%%%%%%%%%%%%%%%%%%%%%%%%%%%%%%%%%%%%%%%%%%%%%
\subsection{StMcCalorimeterHit}
\index{Calorimeter hit} \index{StMcCalorimeterHit|textbf}
\label{sec:StMcCalorimeterHit}
\begin{Entry}
\item[Summary]
    \name{StMcCalorimeterHit} represents a Calorimeter hit.

\item[Synopsis]
    \verb+#include "StMcCalorimeterHit.hh"+\\
    \verb+class StMcCalorimeterHit;+\\

\item[Description]
    \name{StMcCalorimeterHit} does not inherit from \name{StMcHit}.  This decision was
    taken by the EMC group based on the different functionality the Calorimeter Hit
    would need to have.  For example, one of the main difference is that the calorimeter hits
    do not have a position in global coordinates.  They do however, still keep a pointer to
    their parent track.

\item[Persistence]
    None

\item[Related Classes]
    The hits are kept according to module.
    They are stored by pointer in a container for each module
    (see \ref{sec:containers}).  The way to get the hits for a
    certain module (0-119, 120 modules in total) starting
    from the top level \name{StMcEvent} pointer is:
    \verb+mcEvt->emcHitCollection()->module(0)->hits()+
    Look in \ref{sec:StMcEmcHitCollection} and
    \ref{sec:StMcEmcModuleHitCollection} for more information.
    Each instance of \name{StMcTrack} holds a list of Calorimeter hits
    which belong to that track.  Depending on which detector they are, they
    will be in the appropriate containter: Bemc, Bprs, Bsmde, or Bsmdp.  Look
    in \ref{sec:StMcTrack} for more information.
    \index{StMcTrack}
    \index{StMcEmcHitCollection}

\item[Public\\ Constructors]
    \verb+StMcCalorimeterHit(int m,int e,int s,float de)+\\
    Create an instance of \name{StMcCalorimeterHit} with module \name{p},
    eta \name{e}, sub \name{s}, and energy deposition at hit \name{de}.

    \verb+StMcCalorimeterHit(int m,int e,int s,float de, StMcTrack* parent)+\\
    Create an instance of \name{StMcCalorimeterHit} with module \name{p},
    eta \name{e}, sub \name{s}, energy deposition at hit \name{de},
    and parent track \name{parent}.

\item[Public Member\\ Functions]

    \verb+int module() const;+\\
    Returns the number of the module (1-120) in which a hit is found. 
    \index{module, EMC}

    \verb+int eta() const;+\\
    Returns the eta in which a hit is found. 
    \index{eta, EMC}

    \verb+int sub() const;+\\
    Returns the sub in which a hit is found. 
    \index{sub, EMC}

    \verb+float dE() const;+\\
    Returns the energy deposition of the hit.
    \index{dE, EMC}

    \verb+StMcTrack*  parentTrack() const;+\\
    Returns the pointer to the track which generated this hit.
    \index{parentTrack, EMC}

\item[Public Member\\ Operators]
    \verb+ostream&  operator<<(ostream& os, const StMcCalorimeterHit\&);+\\
    Allows to write some relevant information directly to an output
    stream.
    \verb+void operator+=(const StMcCalorimeterHit\&);+\\
    Allows to incrementally add energy to a calorimeter hit from another
    calorimeter hit.
    \verb+void operator==(const StMcCalorimeterHit\&);+\\
    Compares two calorimeter hits, they're equal when the module, eta, sub
    and parent track pointers are found to be equal.

\item[Examples]
{\footnotesize
\begin{verbatim}
//
//  In this example, we use the methods of the class
//  to access the data.  For concreteness, we use them to fill
//  a Root histogram. 
//  
void
PositionOfHits(StMcTrack *track)
{
   StPtrVecMcCalorimeterHit hits = track->bemcHits();

   StMcCalorimeterHitIterator i;
   StMcCalorimeterHit* currentHit;
   TH2F* myHist = new TH2F("coords. MC","eta vs mod pos. of Hits",120, 1, 121, 100, -2, 2)
   for (i = hits.begin(); i != hits.end(); i++) {
      currentHit = (*i);
      myHist->Fill(currentHit->module(), currentHit->eta());
   }
}
\end{verbatim}
}%footnotesize
\end{Entry}
\clearpage
%%%%%%%%%%%%%%%%%%%%%%%%%%%%%%%%%%%%%%%%%%%%%%%%%%%%%%%%%%%%%%%%%%%%
%
%    Reference: StMcCtbHit
%
%%%%%%%%%%%%%%%%%%%%%%%%%%%%%%%%%%%%%%%%%%%%%%%%%%%%%%%%%%%%%%%%%%%%
\subsection{StMcCtbHit}
\index{StMcCtbHit|textbf}
\label{sec:StMcCtbHit}
\begin{Entry}
\item[Summary]
    \name{StMcCtbHit} represents a CTB hit.
\item[Synopsis]
    \verb+#include "StMcCtbHit.hh"+\\
    \verb+class StMcCtbHit;+\\

\item[Description]
Like most of the hit classes, it represents the Geant hit its related detector, in
this case the CTB.
\item[Persistence]
    None

\item[Related Classes]
The class to hold these hits is \name{StMcCtbHitCollection}, which holds them in a flat
array, no further hierarchy.
\item[Public\\ Constructors]
    \verb+StMcCtbHit(const StThreeVectorF& x,const StThreeVectorF& p, const float de, const float ds,  const long key, const long id, StMcTrack* parent)+\\
    Create an instance of \name{StMcCtbHit} with position \name{x},
    local momentum \name{p}, energy deposited \name{de}, pathlength \name{ds},
    key and id for proper assignment of pointers and the pointer to the partent track..

    \verb+StMcCtbHit(g2t_ctf_hit_st*)+\\
    Create an instance of \name{StMcCtbHit} from the geant tables, this is the
    standard way in \name{StMcEventMaker}.

\item[Public Member\\ Functions]
    \verb+void get_slat_tray(unsigned int & slat, unsigned int & tray) const;+\\
    Puts the corresponding slat and tray of the hit in the variables entered as arguments to the function.

    \verb+float tof() const;+\\
    Returns the Time of Flight of the particle to the detector hit volume.

\end{Entry}
%%%%%%%%%%%%%%%%%%%%%%%%%%%%%%%%%%%%%%%%%%%%%%%%%%%%%%%%%%%%%%%%%%%%
%
%    Reference: className  [TEMPLATE PLACEHOLDER — suppressed]
%
%    This was an unfilled template left in the original document.
%    New class entries should follow the \subsection{ClassName} structure
%    shown throughout this section.
%
%%%%%%%%%%%%%%%%%%%%%%%%%%%%%%%%%%%%%%%%%%%%%%%%%%%%%%%%%%%%%%%%%%%%
\subsection{StMcEmcHitCollection}
\index{EMC hit collection} \index{StMcEmcHitCollection|textbf}
\label{sec:StMcEmcHitCollection}
\begin{Entry}
\item[Summary]
    \name{StMcEmcHitCollection} is a container for the
    modules of the EMC.  The actual hits are stored according to
    the module they belong to.

\item[Synopsis]
    \verb+#include "StMcEmcHitCollection.hh"+\\
    \verb+class StMcEmcHitCollection;+\\

\item[Description]
    \name{StMcEmcHitCollection} is the first step down
    the hierarchy of hits for the EMC.  It provides
    methods to return a particular module, and to
    count the number of hits in the EMC detector that it
    represents.  The EMC is different than the rest of the
    detectors in that the same class is used for the Barrel Emc,
    the Pre-shower, and the Shower Max detectors (Eta and Phi).

\item[Persistence]
    None

\item[Related Classes]
    \name{StMcEmcHitCollection}
    has a container of type {\tt StMcEmcModuleHitCollection}
    of size 120, each element represents a module.  Look in
    \name{StMcEmcModuleHitCollection} \ref{sec:StMcEmcModuleHitCollection}.
    This class inherits from \name{TDataSet}, so that each hit collection
    also has a name to distinguish them (along with other inherited properties
    of \name{TDataSet}.
    \index{StMcEmcModuleHitCollection}
    \index{StMcCalorimeterHit}

\item[Public\\ Constructors]
    \verb+StMcTpcHitCollection();+\\
    Create an instance of \name{StEmcHitCollection}
    and sets up the internal containers.

    \verb+StMcTpcHitCollection(char* name);+\\
    Create an instance of \name{StEmcHitCollection}, where \name{name} is passed to the
    TDataSet constructor, 
    and sets up the internal containers.

\item[Enumerations]
    \verb+enum  EAddHit {kNull, kErr, kNew, kAdd};+\\
    Enumerations for the return values of the addHit method to flag
    the different cases.  See below.
\item[Public Member\\ Functions]

    \verb+StMcEmcHitCollection::EAddHit addHit(StMcTpcHit*);+\\
    Adds a hit to the collection (using \name{StMemoryPool}
    from the StarClassLibrary.  If the hit is a new hit, kNew is returned.
    If the hit was already present and we just needed to add the information,
    kAdd is returned.  If the hit has a wrong module or some other parameter that
    prevents its correct assignment, kErr is returned.

    \verb+unsigned long numberOfHits() const;+\\
    Counts the number of hits of the for the referring EMC detector.

    \verb+unsigned int  numberOfModule() const;+\\
    Returns the size of the module vector.  Default is 120.

    \verb+StMcEmcModuleHitCollection*       module(unsigned int);+\\
    Returns a pointer to the module specified by the argument to the
    function.  Recall that this is for indexing into an array, so the
    argument should go from 0 - (size()-1).  Look in the
    numbering scheme conventions \ref{sec:numberscheme}for more information.

\item[Examples]
{\footnotesize
\begin{verbatim}
//
// Look in StMcTpcModuleHitCollection for an example.
// These classes normally go hand in hand.
\end{verbatim}
}%footnotesize
\end{Entry}

\clearpage
%%%%%%%%%%%%%%%%%%%%%%%%%%%%%%%%%%%%%%%%%%%%%%%%%%%%%%%%%%%%%%%%%%%%
%
%    Reference: StMcEmcModuleHitCollection
%
%%%%%%%%%%%%%%%%%%%%%%%%%%%%%%%%%%%%%%%%%%%%%%%%%%%%%%%%%%%%%%%%%%%%
\subsection{StMcEmcModuleHitCollection}
\index{EMC module hit collection} \index{StMcEmcModuleHitCollection|textbf}
\label{sec:StMcEmcModuleHitCollection}
\begin{Entry}
\item[Summary]
    \name{StMcEmcModuleHitCollection} is the class where the EMC hits
    are actually stored.  It holds a vector of StMcCalorimeterHit pointers.

\item[Synopsis]
    \verb+#include "StMcEmcModuleHitCollection.hh"+\\
    \verb+class StMcEmcModuleHitCollection;+\\

\item[Description]
    \name{StMcEmcModuleHitCollection} holds the EMC hits.

\item[Persistence]
    None

\item[Related Classes]
    \name{StMcEmcModuleHitCollection}
    has a container of type {\tt StSPtrVecMcCalorimeterHit}, this means
    that this is the class that actually owns the hits.  Look
    at the section on containers \ref{sec:containers} for more
    information.    
    \index{StMcCalorimeterHit}
    \index{StSPtrVecMcCalorimeterHit}
\item[Public\\ Constructors]
    \verb+StMcEmcModuleHitCollection();+\\
    Create an instance of \name{StEmcHitCollection}
    and sets up the internal container.

\item[Public Member\\ Functions]

    \verb+unsigned long numberOfHits() const;+\\
    Counts the number of hits of this module.

    \verb+float sum() const;+\\
    Returns the summation of the energy deposition, \name{dE}, values
    of all the hits in the module.

    \verb+StSPtrVecMcEmcHit&       hits();+\\
    Returns a reference to the container of the EMC hits.

    
\item[Examples]
{\footnotesize
\begin{verbatim}
//
// Print number of hits in EMC and plot 1D Histogram of eta of
// Hits
//  
   StMcEmcHitCollection* emcHitColl = evt->bemcHits();
   cout << ``There are :`` << emcHitColl->numberOfHits() << ``hits in the BEMC'' << endl; 
   TH1F* etaHist = new TH1F("etaHist","Eta pos. of Hits",100, -2, 2)
   for (unsigned int i=0; i<emcHitColl->numberOfModules(); i++) {
     StSPtrVecMcEmcHit& hits = emcHitColl->module(i)->hits();
     for (StMcEmcHitIterator i = hits.begin(); i != hits.end(); i++) {
       currentHit = (*i);
       myHist->Fill(currentHit->eta());
     }
   }
\end{verbatim}
}%footnotesize
\end{Entry}

\clearpage


%%%%%%%%%%%%%%%%%%%%%%%%%%%%%%%%%%%%%%%%%%%%%%%%%%%%%%%%%%%%%%%%%%%%
%
%    Reference: StMcEvent
%
%%%%%%%%%%%%%%%%%%%%%%%%%%%%%%%%%%%%%%%%%%%%%%%%%%%%%%%%%%%%%%%%%%%%
\subsection{StMcEvent}
\index{StMcEvent|textbf}
\index{event header}
\label{sec:StMcEvent}
\begin{Entry}
\item[Summary]
    \name{StMcEvent} is the top class in the \StMcEvent\ data model.
    It provides methods to access all quantities and objects
    stored in the g2t tables.

\item[Synopsis]
    \verb+#include "StMcEvent.hh"+\\
    \verb+class StMcEvent;+\\

\item[Description]
    Objects of type \name{StMcEvent} are the entry point to the g2t data.
    From here one can navigate to (and access) quantities stored
    in the g2t tables, but instead of going through foreign keys and
     Id's all relationships are established through pointers. 
    The only {\tt \#include} file needed to access all the classes of \StMcEvent\
    is the {\tt StMcEventTypes.hh} file. \index{StMcEventTypes}
    \name{StMcEvent} itself doesn't offer much functionality but rather serves
    as a container for the event data.
    Deleting \name{StMcEvent} deletes the whole data tree, i.e.~all depending objects
    are deleted properly.

\item[Persistence]
    None

\item[Related Classes]
    None

\item[Public\\ Constructors]
    \verb+StMcEvent();+\\
    Default constructor. Creates an "empty" event.  Normally not used.  As with
    most other classes, the constructor used normally is based on g2t tables.

\item[Public Member\\ Functions]
    

    \verb+unsigned long eventGeneratorEventLabel() const;+\\
    Event Label from event generator, read from g2t\_event.
    \index{event generator label}

    \verb+unsigned long eventNumber() const;+\\
    Event number from monte carlo production, read from g2t\_event.
    \index{event number}

    \verb+unsigned long runNumber() const;+\\
    Run number, read from g2t\_event.
    \index{run number}

    \verb+unsigned long zWest() const;+\\
    Number of protons of particle coming from the ``West'', read from g2t\_event.
    \index{zWest}

    \verb+unsigned long nWest() const;+\\
    Number of neutrons of particle coming from the ``West'', read from g2t\_event.
    \index{nWest}

    \verb+unsigned long zEast() const;+\\
    Number of protons of particle coming from the ``East'', read from g2t\_event.
    \index{zEast}

    \verb+unsigned long nEast() const;+\\
    Number of neutrons of particle coming from the ``East'', read from g2t\_event.
    \index{nEast}

    \verb+unsigned long eventGeneratorFinalStateTracks() const;+\\
    As the name implies, this has the number of final stated tracks coming
    from the event generator, read from g2t\_event.
    \index{event generator final state tracks}

    \verb+unsigned long numberOfPrimaryTracks() const;+\\
    Number of primary tracks read from g2t\_event.  Note that this is
    not coming from {\tt primaryVertex()->daughters().size()} so one
    can check whether these numbers are actually the same (they should be).  Note
    that the tracks that only appear in the particle table are NOT assigned
    either a start or a stop vertex, and they are NOT kept as daughters
    of the primary vertex either.
    \index{primary tracks}

    \verb+unsigned long subProcessId() const;+\\
    This is used for simulated Pythia pp events, where
    one can study a given reaction, e.g. gg->QQbar = 82, and
    gg->J/Psi+g = 86. Refer to the Pythia manual for the
    subprocess ids.  Read from g2t\_event.
    \index{subprocess ID}
    \index{Pythia}

    \verb+float impactParameter() const;+\\
    Impact parameter of collision, read from g2t\_event.
    \index{impact parameter}

    \verb+float phiReactionPlane() const;+\\
    Phi angle of the reaction plane of the collision, read from g2t\_event.
    \index{phi reaction plane}

    \verb+float triggerTimeOffset() const;+\\
    Trigger time offset, read from g2t\_event.
    \index{trigger time offset}

    \verb+unsigned long nBinary() const;+\\
    Returns the number of binary collisions.  Read from g2t\_event.
    \index{binary collisions}

    \verb+unsigned long nWoundedEast() const;+\\
    \verb+unsigned long nWoundedWest() const;+\\
    Return the number of binary of wounded nucleons from each collision
    direction.  Read from g2t\_event.
    \index{wounded nucleons}

    \verb+unsigned long nJets() const;+\\
    Number of jets in the event.  Read from g2t\_event.
    \index{jets}



    \verb+StMcVertex* primaryVertex();+\\
    Returns a pointer to primary vertex. The same object is stored
    in the \name{StMcVertexCollection}. Since the primary vertex is of
    high importance for most analysis steps, this method was added
    for convenience. Note that if no primary vertex is available
    this function returns a {\tt NULL} pointer, so be careful!
    \index{primary vertex}

    \verb+StSPtrVecMcVertex& vertices();+\\
    Returns a reference to the vertex container, i.e., the {\tt vector} of all
    event vertices.
    \index{vertex container}  
    \verb+StSPtrVecMcTrack& tracks();+\\
    Returns reference to the track container, i.e., the {\tt vector} of all
    monte carlo tracks.
    \index{track container}

    \verb+StMcTpcHitCollection* tpcHitCollection();+\\ 
    Returns a pointer to the TPC hit collection.  Note that the
    collections are {\emph NOT} just flat containers.  The TPC hit
    collection is a container of sectors.  The sectors are containers
    of padrows, where the TPC hits are stored.
    \index{hit collections, TPC}

    \verb+StMcSvtHitCollection* svtHitCollection();+\\
    Returns a pointer to the SVT hit collection.  The SVT hit
    collection is a container of layers.  The layers are containers
    of ladders.  The ladders are containers of wafers,
    where the SVT hits are stored.
    \index{hit collections, SVT}


    \verb+StMcFtpcHitCollection* ftpcHitCollection();+\\
    Returns a pointer to the FTPC hit collection. The FTPC hit
    collection is a container of planes.  Note that it is in the planes 
    where the FTPC hits are stored.  The Monte Carlo hits of the FTPC
    know nothing about sectors, because these are only determined
    after reconstruction.
    \index{hit collections, FTPC}

    \verb+StMcRichHitCollection* richHitCollection();+\\
    Returns a pointer to the RICH hit collection. This is a flat collection
    of hits.

    \verb+StMcEmcHitCollection* bemcHitCollection();+\\
    Returns a pointer to the BEMC hit collection. The calorimeter hit
    collection is containers of modules.  This same class is used for
    the other calorimeter detector components, below.

    \verb+StMcEmcHitCollection* bprsHitCollection();+\\
    Returns a pointer to the Barrel Pre-shower hit collection.

    \verb+StMcEmcHitCollection* bsmdeHitCollection();+\\
    Returns a pointer to the Barrel Shower Max detector-Eta hit collection.

    \verb+StMcEmcHitCollection* bsmdpHitCollection();+\\
    Returns a pointer to the Barrel Shower Max detector-Phi hit collection.

    The following member functions are used to set data members of \name{StMcEvent}.
    They are shown only for completeness and shouldn't be used without
    a profound understanding of the relations between the different objects.

    \verb+void setEventGeneratorEventLabel(unsigned long);+\\
    \verb+void setEventNumber(unsigned long);+\\
    \verb+void setRunNumber(unsigned long);+\\
    \verb+void setZWest(unsigned long);+\\
    \verb+void setNWest(unsigned long);+\\
    \verb+void setZEast(unsigned long);+\\
    \verb+void setNEast(unsigned long);+\\
    \verb+void setImpactParameter(float);+\\
    \verb+void setPhiReactionPlane(float);+\\
    \verb+void setTriggerTimeOffset(float);+\\
    \verb+void setPrimaryVertex(StMcVertex*);+\\
    \verb+void setTrackCollection(StMcTrackCollection*);+\\
    \verb+void setTpcHitCollection(StMcTpcHitCollection*);+\\
    \verb+void setSvtHitCollection(StMcSvtHitCollection*);+\\
    \verb+void setFtpcHitCollection(StMcFtpcHitCollection*);+\\
    \verb+void setRichHitCollection(setRichHitCollection);+\\
    \verb+void setBemcHitCollection(StMcEmcHitCollection);+\\
    \verb+void setBprsHitCollection(StMcEmcHitCollection);+\\
    \verb+void setBsmdeHitCollection(StMcEmcHitCollection);+\\
    \verb+void setBsmdpHitCollection(StMcEmcHitCollection);+\\
    \verb+void setVertexCollection(StMcVertexCollection*);+\\

\item[Public Member\\ Operators]
    \verb+int operator==(const StMcEvent&) const;+\\
    Compares the event number, run number and type to determine
    whether 2 events are equal.

    \verb+int operator!=(const StMcEvent&) const;+\\
    Uses {\tt operator==} and negates it.

    \verb+ostream&  operator<<(ostream& os, const StMcEvent&);+\\
    Allows to write some relevant information directly to an output
    stream.

\item[Examples]
{\bf Example:}
{\footnotesize
\begin{verbatim}
//
//   Prints some basic event quantities to stdout (cout).
//   
//
void printEvent(StMcEvent *event)
{
     
     cout << "Event Generator Label = "
          <<  event->eventGeneratorEventLabel() << endl;
     cout << "Event Number = "
          <<  event->eventNumber() << endl;
     cout << "Run number = " << event->runNumber();
     cout << "# of tracks = "
          << event->trackCollection()->size() << endl;
     cout << "# of vertices = "
          << event->vertexCollection()->size() << endl;
     cout << "# of tpc hits = "
          << event->tpcHitCollection()->size() << endl;
     cout << "xyz of primary Vertex = "
          <<  event->primaryVertex()->position()/millimeter
          << " mm" << endl;
}

\end{verbatim}
}%footnotesize


\end{Entry}
\clearpage

%%%%%%%%%%%%%%%%%%%%%%%%%%%%%%%%%%%%%%%%%%%%%%%%%%%%%%%%%%%%%%%%%%%%
%
%    Reference: StMcFtpcHit
%
%%%%%%%%%%%%%%%%%%%%%%%%%%%%%%%%%%%%%%%%%%%%%%%%%%%%%%%%%%%%%%%%%%%%
\subsection{StMcFtpcHit}
\index{FTPC hit} \index{StMcFtpcHit|textbf}
\label{sec:StMcFtpcHit}
\begin{Entry}
\item[Summary]
    \name{StMcFtpcHit} represents a FTPC hit.

\item[Synopsis]
    \verb+#include "StMcFtpcHit.hh"+\\
    \verb+class StMcFtpcHit;+\\

\item[Description]
    \name{StMcFtpcHit} inherits most functionality from \name{StMcHit}.  For a complete
    description of the inherited member functions see \name{StMcHit}
    (sec.~\ref{sec:StMcHit}).

    In the \StMcEvent\ data model each track keeps references to the
    associated hits and vice versa. All hits have a pointer to their
    parent track.  This is especially important for the making of the
    associations in \StAssociationMaker.


\item[Persistence]
    None

\item[Related Classes]
    \name{StMcFtpcHit} is derived directly from \name{StMcHit}.
    The hits are kept according to plane.
    They are stored by pointer in a container for each plane
    (see \ref{sec:containers}).  The way to get the hits for a
    certain plane (1-20) starting
    from the top level \name{StMcEvent} pointer is:
    \verb+mcEvt->ftpcHitCollection()->plane(0)->hits()+
    Look in \ref{sec:StMcFtpcHitCollection} and
    \ref{sec:StMcFtpcPlaneHitCollection} for more information.
    Each instance of \name{StMcTrack} holds a list of FTPC hits
    which belong to that track.
    \index{StMcHit}
    \index{StMcTrack}
    \index{StMcFtpcHitCollection}

\item[Public\\ Constructors]
    \verb+StMcFtpcHit(const StThreeVectorF& x,const StThreeVectorF& p, +\\
    \verb+          const float de, const float ds, const long key,+\\
    \verb+          const long id, StMcTrack* parent);+\\
    Create an instance of \name{StFtpcHit} with position \name{x}, local momentum \name{p},
    energy deposition at hit \name{de}, path length (within padrow) \name{ds}, primary key \name{key},
    volume Id \name{id},
    and parent track \name{parent}.  Not used.  Standard use is
    to obtain all parameters from g2t\_fpt\_hit table.

\item[Public Member\\ Functions]

    \verb+unsigned long plane() const;+\\
    Returns the number of the plane (1-10 for West Ftpc, 11-20 for East Ftpc) in which a hit is found. 
    \index{plane, FTPC}
    \verb+unsigend long sector() const;+\\
    Returns the sector (1-6) in which a hit is found.  The volume Id to allow this was introduced on Oct. 2003.

\item[Public Member\\ Operators]
    \verb+ostream&  operator<<(ostream& os, const StMcFtpcHit\&);+\\
    Allows to write some relevant information directly to an output
    stream.

\item[Examples]
{\footnotesize
\begin{verbatim}
//
//  In this example, we use the methods of the class
//  to access the data.  For concreteness, we use them to fill
//  a Root histogram. 
//  
void
PositionOfHits(StMcTrack *track)
{
   StPtrVecMcFtpcHit hits = track->ftpcHits();

   StMcFtpcHitIterator i;
   StMcFtpcHit* currentHit;
   TH2F* myHist = new TH2F("coords. MC","X vs Y pos. of Hits",100, -150, 150, 100, -150, 150)
   for (i = hits.begin(); i != hits.end(); i++) {
      currentHit = (*i);
      myHist->Fill(currentHit->position().x(), currentHit->position().y());
   }
}
\end{verbatim}
}%footnotesize
\end{Entry}

\clearpage
%%%%%%%%%%%%%%%%%%%%%%%%%%%%%%%%%%%%%%%%%%%%%%%%%%%%%%%%%%%%%%%%%%%%
%
%    Reference: StMcFtpcHitCollection
%
%%%%%%%%%%%%%%%%%%%%%%%%%%%%%%%%%%%%%%%%%%%%%%%%%%%%%%%%%%%%%%%%%%%%
\subsection{StMcFtpcHitCollection}
\index{FTPC hit collection} \index{StMcFtpcHitCollection|textbf}
\label{sec:StMcFtpcHitCollection}
\begin{Entry}
\item[Summary]
    \name{StMcFtpcHitCollection} is a container for the
    planes of the FTPC.  The actual hits are stored according to
    planes.

\item[Synopsis]
    \verb+#include "StMcFtpcHitCollection.hh"+\\
    \verb+class StMcFtpcHitCollection;+\\

\item[Description]
    \name{StMcFtpcHitCollection} is the first step down
    the hierarchy of hits for the FTPC.  It provides
    methods to return a particular plane, and to
    count the number of hits in all the FTPC.

\item[Persistence]
    None

\item[Related Classes]
    \name{StMcFtpcHitCollection}
    has a container of type {\tt StMcFtpcPlaneHitCollection}
    of size 20, each element represents a plane.
    \index{StMcFtpcPlaneHitCollection}
    \index{StMcFtpcHit}

\item[Public\\ Constructors]
    \verb+StMcFtpcHitHitCollection();+\\
    Create an instance of \name{StFtpcHitCollection}
    and sets up the internal containers.

\item[Public Member\\ Functions]

    \verb+bool addHit(StMcFtpcHit*);+\\
    Adds a hit to the collection (using \name{StMemoryPool}
    from the StarClassLibrary.

    \verb+unsigned long numberOfHits() const;+\\
    Counts the number of hits of the whole FTPC.

    \verb+unsigned int  numberOfPlanes() const;+\\
    Returns the size of the Array of planes.  Default is 20.

    \verb+StMcFtpcPlaneHitCollection*       plane(unsigned int);+\\
    Returns a pointer to the plane specified by the argument to the
    function.  Recall that this is for indexing into an array, so the
    argument should go from 0 - (size()-1).  Look in the
    numbering scheme conventions \ref{sec:numberscheme}for more information.

\item[Examples]
{\footnotesize
\begin{verbatim}
//
// Look in StMcFtpcPlaneHitCollection for an example.
// These classes normally go hand in hand.
\end{verbatim}
}%footnotesize
\end{Entry}

\clearpage
%%%%%%%%%%%%%%%%%%%%%%%%%%%%%%%%%%%%%%%%%%%%%%%%%%%%%%%%%%%%%%%%%%%%
%
%    Reference: StMcFtpcPlaneHitCollection
%
%%%%%%%%%%%%%%%%%%%%%%%%%%%%%%%%%%%%%%%%%%%%%%%%%%%%%%%%%%%%%%%%%%%%
\subsection{StMcFtpcPlaneHitCollection}
\index{FTPC plane hit collection} \index{StMcFtpcPlaneHitCollection|textbf}
\label{sec:StMcFtpcPlaneHitCollection}
\begin{Entry}
\item[Summary]
    \name{StMcFtpcPlaneHitCollection} is the class where the FTPC hits
    are actually stored.

\item[Synopsis]
    \verb+#include "StMcFtpcPlaneHitCollection.hh"+\\
    \verb+class StMcFtpcPlaneHitCollection;+\\

\item[Description]
    \name{StMcFtpcPlaneHitCollection} holds the FTPC hits.  Note
    the difference with \StEvent.  The Monte Carlo doesn't know
    about the sectors of the FTPC, these are determined during
    reconstruction.

\item[Persistence]
    None

\item[Related Classes]
    \name{StMcFtpcPlaneHitCollection}
    has a container of type {\tt StSPtrVecMcFtpcHit}, this means
    that this is the class that actually owns the hits.  Look
    at the section on containers \ref{sec:containers} for more
    information.    
    \index{StMcFtpcHit}
    \index{StSPtrVecMcFtpcHit}
\item[Public\\ Constructors]
    \verb+StMcFtpcPlaneHitCollection();+\\
    Create an instance of \name{StFtpcHitCollection}
    and sets up the internal container.

\item[Public Member\\ Functions]

    \verb+unsigned long numberOfHits() const;+\\
    Counts the number of hits of this plane.

    \verb+StSPtrVecMcFtpcHit&       hits();+\\
    Returns a reference to the container of the FTPC hits.
    
\item[Examples]
{\footnotesize
\begin{verbatim}
//
// Print number of hits in FTPC and plot 1D Histogram of r (cylindrical coord.) of
// Hits
//  
   StMcFtpcHitCollection* ftpcHitColl = evt->ftpcHits();
   cout << ``There are :`` << ftpcHitColl->numberOfHits() << ``hits in the FTPC'' << endl; 
   TH1F* myHist = new TH1F("Sector 10 ","R pos. of Hits",100, 0, 100)
   for (unsigned int i=0; i<ftpcHitColl->numberOfPlanes(); i++) {
     StSPtrVecMcFtpcHit& hits = ftpcHitColl->plane(i)->hits();
     for (StMcFtpcHitIterator i = hits.begin(); i != hits.end(); i++) {
       currentHit = (*i);
       myHist->Fill(currentHit->position().perp());
     }
   }
\end{verbatim}
}%footnotesize
\end{Entry}

\clearpage
%%%%%%%%%%%%%%%%%%%%%%%%%%%%%%%%%%%%%%%%%%%%%%%%%%%%%%%%%%%%%%%%%%%%
%
%    Reference: StMcHit
%
%%%%%%%%%%%%%%%%%%%%%%%%%%%%%%%%%%%%%%%%%%%%%%%%%%%%%%%%%%%%%%%%%%%%
\subsection{StMcHit}
\index{StMcHit|textbf} \index{hit}
\label{sec:StMcHit}
\begin{Entry}
\item[Summary]
    \name{StMcHit} is the base class of all TPC, FTPC, SVT and RICH
    monte carlo hit classes.

\item[Synopsis]
    \verb+#include "StMcHit.hh"+\\
    \verb+class StMcHit;+\\

\item[Description]
    \name{StMcHit} provides the basic functionality for all derived classes
    which represent TPC, FTPC SVT and RICH hits. It provides information on
    position, energy deposition, path length within sensitive volume,
    and the track which
    generated this hit.  It also has the primary key from the g2t tables and the
    volume Id, which are useful for debugging purposes. All information is taken from
    the appropriate \name{g2t\_xxx\_hit.idl} table, where xxx can be ``tpc'', ``svt'',
    ``ftp'', or ``rch''.
    \name{StMcHit} is a virtual class.
    \index{g2t\_xxx\_hit.idl}

\item[Persistence]
    None

\item[Related Classes]
    The following classes are derived from \name{StMcHit}:
    \begin{itemize}
    \item \name{StMcTpcHit}
    \item \name{StMcFtpcHit}
    \item \name{StMcSvtHit}
    \item \name{StMcRichHit}
    \end{itemize}
    \index{StMcTpcHit}
    \index{StMcFtpcHit}
    \index{StMcSvtHit}
    \index{StMcRichHit}

\item[Public\\ Constructors]
    \verb+StMcHit();+\\
    Constructs an instance of \name{StMcHit} with all values initialized
    to 0 (zero).  Not used.  Normally, all hits belong to a certain detector,
    so the appropriate table based constructor is used.

    \verb+StMcHit(const StThreeVectorF& x, const StThreeVectorF& p,+\\
    \verb+      float de, float ds,  long k, long volId, StMcTrack* parent);+\\
    
    Create an instance of \name{StMcHit} with position \name{x}, local momentum \name{p},
    energy deposition \name{de}, path length \name{ds}, primary key \name{k},
    volume id \name{volId} and parent track \name{parent}.  Used by the derived
    classes to initialize the base class data members.

\item[Public Member\\ Functions]
    \verb+const StThreeVectorF& position() const;+\\
    Hit position in global coordinates.

    \verb+const StThreeVectorF& localMomentum() const;+\\
    Local momentum direction of the parent track at the position of the hit.

    \verb+float dE() const;+\\
    Energy deposition.

    \verb+float dS() const;+\\
    Path length within sensitive volume.

    \verb+long key() const;+\\
    Primary key from the appropriate g2t table.

    \verb+long volumeId() const;+\\
    Geant volume Id, without any decoding.

    \verb+StMcTrack* parentTrack() const;+\\
    Pointer to the track that generated the hit.

    The following functions should not normally be used.
    The relevant data members are set during construction or
    in \name{StMcEventMaker} so these are provided for completeness.

    \verb+void setPosition(const StThreeVectorF&);+\\
    Set the hit position.

    \verb+void setLocalMomentum(const StThreeVectorF&);+\\
    Set the local momentum.

    \verb+void setdE(float);+\\
    Set the total energy deposition.

    \verb+void setdS(float);+\\
    Set the path length within sensitive volume.

    \verb+void setKey(long);+\\
    Set the primary key.

    \verb+void setVolumeId(long);+\\
    Set the volume Id.

    \verb+void setParentTrack(StMcTrack*);+\\
    Set the parent track.
    
\item[Global Operators]

    \verb+int operator==(const StMcHit&) const;+\\
    Checks to see if 2 hits are equal.  Uses position, dE and dS
    for the comparison.

    \verb+int operator!=(const StMcHit&) const;+\\
    Uses operator== and negates the result.

    \verb+ostream&  operator<<(ostream& os, const StMcHit&);+\\
    Prints hit data (position, dE, dS ) to output stream \name{os}.

\item[Examples]
{\footnotesize
\begin{verbatim}
//
//  Create a random hit and print it.
//
void printArbitraryHit()
{
   HepJamesRandom engine;
   RandFlat       rndm(engine);
   RandGauss      rgauss(engine);
   const double   sigmadE = 1e-06;
   const double   sigmadS = 7*millimeter;
   StMcTrack* pTrack = new StMcTrack;

   StThreeVectorF pos(rndm.shoot(-meter,meter),
                            rndm.shoot(-meter,meter),
                            rndm.shoot(-meter,meter));
   StThreeVectorF mom(rndm.shoot(-GeV,GeV),
                            rndm.shoot(-GeV,GeV),
                            rndm.shoot(-GeV,GeV));
   float de = rgauss.shoot(2.0e-06,sigmadE);
   float ds = rgauss.shoot(1.6,sigmadS);
   StMcHit hit(pos, mom, de, ds, 0, 0, pTrack);

   cout << hit << endl;
}
\end{verbatim}
}%footnotesize
{\bf Programs Output:}
{\footnotesize
\begin{verbatim}
Id: 0
Position: -20.3111  -80.7543 76.8227
Local Momentum: .9042 -.4982 -.1873
dE: 1.0952e-06
dS: 1.3049
Vol Id: 0
\end{verbatim}
}%footnotesize

\end{Entry}
\clearpage
%%%%%%%%%%%%%%%%%%%%%%%%%%%%%%%%%%%%%%%%%%%%%%%%%%%%%%%%%%%%%%%%%%%%
%
%    Reference: StMcRichHit
%
%%%%%%%%%%%%%%%%%%%%%%%%%%%%%%%%%%%%%%%%%%%%%%%%%%%%%%%%%%%%%%%%%%%%
\subsection{StMcRichHit}
\index{RICH hit} \index{StMcRichHit|textbf}
\label{sec:StMcRichHit}
\begin{Entry}
\item[Summary]
    \name{StMcRichHit} represents a monte carlo RICH hit.

\item[Synopsis]
    \verb+#include "StMcRichHit.hh"+\\
    \verb+class StMcRichHit;+\\

\item[Description]    
    \name{StMcRichHit} inherits most functionality from \name{StMcHit}.
    For a complete description of the inherited member functions see
    \name{StMcHit} (sec.~\ref{sec:StMcHit}).

    In the \StMcEvent\ data model each track keeps references to the
    associated hits and vice versa. All hits have a pointer to their
    parent track.  This is especially important for the making of the
    associations in \StAssociationMaker. 
    
    
    
\item[Persistence]
    None

\item[Related Classes]
    \name{StMcRichHit} is derived directly from \name{StMcHit}.
    The hits are kept according to row and pad.
    The hits are kept in an instance of \name{StSPtrVecMcRichHit}
    where they are stored by pointer (see \ref{sec:containers}).
    Each instance of \name{StMcTrack} holds a list of RICH hits
    which belong to that track.
    \index{StMcHit}
    \index{StMcTrack}
    \index{StMcRichHitCollection}

\item[Public\\ Constructors]
    \verb+StMcRichHit(const StThreeVectorF& x,const StThreeVectorF& p,+\\
    \verb+      const float de, const float ds, const long key, const long id, StMcTrack* parent);+\\
    
    Create an instance of \name{StMcHit} with position \name{x}, local momentum \name{p},
    energy deposition \name{de}, path length \name{ds}, primary key \name{key},
    volume id \name{id} and parent track \name{parent}.  

\item[Public Member\\ Functions]

    \verb+unsigned short  pad() const;+\\
    Returns the number of the pad in which a hit is found. 
    \index{pad, RICH Hit}

    \verb+unsigned short  row() const;+\\
    Returns the number of the row in which a hit is found. 
    \index{row, RICH Hit}

    \verb+float  tof() const;+\\
    Returns the time of flight from the g2t\_rch\_hit table.
    \index{tof, RICH Hit}

\item[Examples]
{\footnotesize
\begin{verbatim}
//
//  Function which returns the RICH hits
//  used by a track and prints them.
//
void getHits(StMcTrack *track)
{
    StPtrVecMcRichHit& hits =
                   track->richHits();
    StMcRichHitIterator iter;
    StMcRichHit* hit;
    for(iter = hits->begin();
        iter != hits->end(); iter++) {
        hit = *iter;
        cout << *hit << endl;
    }
    
}
\end{verbatim}
}%footnotesize
\end{Entry}
\clearpage
%%%%%%%%%%%%%%%%%%%%%%%%%%%%%%%%%%%%%%%%%%%%%%%%%%%%%%%%%%%%%%%%%%%%
%
%    Reference: StMcRichHitCollection
%
%%%%%%%%%%%%%%%%%%%%%%%%%%%%%%%%%%%%%%%%%%%%%%%%%%%%%%%%%%%%%%%%%%%%
\subsection{StMcRichHitCollection}
\index{RICH hit collection} \index{StMcRichHitCollection|textbf}
\label{sec:StMcRichHitCollection}
\begin{Entry}
\item[Summary]
    \name{StMcRichHitCollection} is a container for the
    hits of the RICH.  There is no hierarchy, this class holds
    the vector of hits.

\item[Synopsis]
    \verb+#include "StMcRichHitCollection.hh"+\\
    \verb+class StMcRichHitCollection;+\\

\item[Description]
    \name{StMcRichHitCollection} holds the hits and the methods to add and count them.

\item[Persistence]
    None

\item[Related Classes]
    \name{StMcRichHitCollection}
    has a container of type {\tt StSPtrVecMcRichHit} which holds the hits.
    \index{StMcRichHit}

\item[Public\\ Constructors]
    \verb+StMcRichHitCollection();+\\
    Create an instance of \name{StRichHitCollection}
    and sets up the internal container.

\item[Public Member\\ Functions]

    \verb+bool addHit(StMcRichHit*);+\\
    Adds a hit to the collection (using \name{StMemoryPool}
    from the StarClassLibrary.

    \verb+unsigned long numberOfHits() const;+\\
    Counts the number of hits of the whole RICH.

    \verb+StSPtrVecMcRichHit& hits();+\\
    Returns a reference to the vector of hits.  This function also has a \verb+const+ version.
    

\item[Examples]
{\footnotesize
\begin{verbatim}
//
// Look in any of the other hit collections, the usage is identical. 
//
\end{verbatim}
}%footnotesize
\end{Entry}


\clearpage
%%%%%%%%%%%%%%%%%%%%%%%%%%%%%%%%%%%%%%%%%%%%%%%%%%%%%%%%%%%%%%%%%%%%
%
%    Reference: StMcSvtHit
%
%%%%%%%%%%%%%%%%%%%%%%%%%%%%%%%%%%%%%%%%%%%%%%%%%%%%%%%%%%%%%%%%%%%%
\subsection{StMcSvtHit}
\index{SVT hit} \index{StMcSvtHit|textbf}
\label{sec:StMcSvtHit}
\begin{Entry}
\item[Summary]
    \name{StMcSvtHit} represents a monte carlo SVT hit.

\item[Synopsis]
    \verb+#include "StMcSvtHit.hh"+\\
    \verb+class StMcSvtHit;+\\

\item[Description]    
    \name{StMcSvtHit} inherits most functionality from \name{StMcHit}.
    For a complete description of the inherited member functions see
    \name{StMcHit} (sec.~\ref{sec:StMcHit}).

    In the \StMcEvent\ data model each track keeps references to the
    associated hits and vice versa. All hits have a pointer to their
    parent track.  This is especially important for the making of the
    associations in \StAssociationMaker. 
    
    
    
\item[Persistence]
    None

\item[Related Classes]
    \name{StMcSvtHit} is derived directly from \name{StMcHit}.
    The way to get the hits for a
    certain layer-ladder-wafer starting
    from the top level \name{StMcEvent} pointer is:
    \verb+mcEvt->svtHitCollection()->layer(ly)->ladder(ld)->wafer(w)->hits()+
    The hits are kept in an instance of \name{StSPtrVecMcSvtHit}
    where they are stored by pointer (see \ref{sec:containers}).
    Each instance of \name{StMcTrack} holds a list of SVT hits
    which belong to that track.
    \index{StMcHit}
    \index{StMcTrack}
    \index{StMcSvtHitCollection}

\item[Public\\ Constructors]
    \verb+StMcSvtHit(const StThreeVectorF& x,const StThreeVectorF& p,+\\
    \verb+      const float de, const float ds, const long key, const long id, StMcTrack* parent);+\\
    
    Create an instance of \name{StMcHit} with position \name{x}, local momentum \name{p},
    energy deposition \name{de}, path length \name{ds}, primary key \name{key},
    volume id \name{id} and parent track \name{parent}. 

\item[Public Member\\ Functions]

    \verb+unsigned long  layer() const;+\\
    Returns the number of the layer [1-6] in which a hit is found. 
    \index{layer, SVT}

    \verb+unsigned long  ladder() const;+\\
    Returns the number of the ladder [1-8] in which a hit is found. 
    \index{ladder, SVT}

    \verb+unsigned long  wafer() const;+\\
    Returns the number of the wafer [1-7] in which a hit is found. 
    \index{wafer, SVT}

    \verb+unsigned long  barrel() const;+\\
    Returns the number of the barrel [1-3] in which a hit is found. 
    \index{barrel, SVT}

    \verb+unsigned long  hybrid() const;+\\
    Returns the number of the hybrid.
    \index{hybrid, SVT}

\item[Examples]
{\footnotesize
\begin{verbatim}
//
//  Function which returns the SVT hits
//  used by a track and prints them.
//
void getHits(StMcTrack *track)
{
    StPtrVecMcSvtHit& hits =
                   track->svtHits();
    StMcSvtHitIterator iter;
    StMcSvtHit* hit;
    for(iter = hits->begin();
        iter != hits->end(); iter++) {
        hit = *iter;
        cout << hit << endl;
    }
    
}
\end{verbatim}
}%footnotesize
\end{Entry}
\clearpage
%%%%%%%%%%%%%%%%%%%%%%%%%%%%%%%%%%%%%%%%%%%%%%%%%%%%%%%%%%%%%%%%%%%%
%
%    Reference: StMcSvtHitCollection
%
%%%%%%%%%%%%%%%%%%%%%%%%%%%%%%%%%%%%%%%%%%%%%%%%%%%%%%%%%%%%%%%%%%%%
\subsection{StMcSvtHitCollection}
\index{SVT hit collection} \index{StMcSvtHitCollection|textbf}
\label{sec:StMcSvtHitCollection}
\begin{Entry}
\item[Summary]
    \name{StMcSvtHitCollection} is a container for the
    layers of the SVT.  The actual hits are stored according to
    wafers, 3 levels down the hierarchy.

\item[Synopsis]
    \verb+#include "StMcSvtHitCollection.hh"+\\
    \verb+class StMcSvtHitCollection;+\\

\item[Description]
    \name{StMcSvtHitCollection} is the first step down
    the hierarchy of hits for the SVT.  It provides
    methods to return a particular layer, and to
    count the number of hits in all the SVT.

\item[Persistence]
    None

\item[Related Classes]
    \name{StMcSvtHitCollection}
    has a container of type {\tt StMcSvtLayerHitCollection}
    of size 6, each element represents a layer.  Look in
    \name{StMcSvtLayerHitCollection} \ref{sec:StMcSvtLayerHitCollection},
    \name{StMcSvtLadderHitCollection} \ref{sec:StMcSvtLadderHitCollection},
    and \name{StMcSvtWaferHitCollection} \ref{sec:StMcSvtWaferHitCollection},
    \index{StMcSvtLayerHitCollection}
    \index{StMcSvtLadderHitCollection}
    \index{StMcSvtWaferHitCollection}
    \index{StMcSvtHit}

\item[Public\\ Constructors]
    \verb+StMcSvtHitCollection();+\\
    Create an instance of \name{StSvtHitCollection}
    and sets up the internal containers.

\item[Public Member\\ Functions]

    \verb+bool addHit(StMcSvtHit*);+\\
    Adds a hit to the collection (using \name{StMemoryPool}
    from the StarClassLibrary.

    \verb+unsigned long numberOfHits() const;+\\
    Counts the number of hits of the whole SVT.

    \verb+unsigned int  numberOfLayers() const;+\\
    Returns the size of the Array of layers.  Default is 6.

    \verb+StMcSvtLayerHitCollection*       layer(unsigned int);+\\
    Returns a pointer to the layer specified by the argument to the
    function.  Recall that this is for indexing into an array, so the
    argument should go from 0 - (size()-1).  Look in the
    numbering scheme conventions \ref{sec:numberscheme}for more information.

\item[Examples]
{\footnotesize
\begin{verbatim}
//
// Look in StMcSvtWaferHitCollection for an example.
// These classes normally go hand in hand.
\end{verbatim}
}%footnotesize
\end{Entry}

\clearpage
%%%%%%%%%%%%%%%%%%%%%%%%%%%%%%%%%%%%%%%%%%%%%%%%%%%%%%%%%%%%%%%%%%%%
%
%    Reference: StMcSvtLadderHitCollection
%
%%%%%%%%%%%%%%%%%%%%%%%%%%%%%%%%%%%%%%%%%%%%%%%%%%%%%%%%%%%%%%%%%%%%
\subsection{StMcSvtLadderHitCollection}
\index{SVT ladder hit collection} \index{StMcSvtLadderHitCollection|textbf}
\label{sec:StMcSvtLadderHitCollection}
\begin{Entry}
\item[Summary]
    \name{StMcSvtLadderHitCollection} is a container for the
    wafers of the SVT.  The actual hits are stored according to
    wafers, the next level down the hierarchy.

\item[Synopsis]
    \verb+#include "StMcSvtLadderHitCollection.hh"+\\
    \verb+class StMcSvtLadderHitCollection;+\\

\item[Description]
    \name{StMcSvtLadderHitCollection} is the third step down
    the hierarchy of hits for the SVT.  It provides
    methods to return a particular wafer, and to
    count the number of hits in the ladder.

\item[Persistence]
    None

\item[Related Classes]
    \name{StMcSvtLadderHitCollection}
    has a container of type {\tt StMcSvtWaferHitCollection}
    of size 7, each element represents a wafer.  Look in
    \name{StMcSvtLayerHitCollection} \ref{sec:StMcSvtLayerHitCollection},
    \name{StMcSvtHitCollection} \ref{sec:StMcSvtHitCollection},
    and \name{StMcSvtWaferHitCollection} \ref{sec:StMcSvtWaferHitCollection},
    \index{StMcSvtLayerHitCollection}
    \index{StMcSvtHitCollection}
    \index{StMcSvtWaferHitCollection}
    \index{StMcSvtHit}

\item[Public\\ Constructors]
    \verb+StMcSvtLadderHitHitCollection();+\\
    Create an instance of \name{StSvtLadderHitCollection}
    and sets up the internal containers.

\item[Public Member\\ Functions]

    \verb+unsigned long numberOfHits() const;+\\
    Counts the number of hits of this ladder.

    \verb+unsigned int  numberOfWafers() const;+\\
    Returns the size of the Array of wafers.  Default is 7.

    \verb+StMcSvtWaferHitCollection*       wafer(unsigned int);+\\
    Returns a pointer to the wafer hit collection
    specified by the argument to the
    function.  Recall that this is for indexing into an array, so the
    argument should go from 0 - (size()-1).  Look in the
    numbering scheme conventions \ref{sec:numberscheme}for more information.

\item[Examples]
{\footnotesize
\begin{verbatim}
//
// Look in StMcSvtWaferHitCollection for an example.
// These classes normally go hand in hand.
\end{verbatim}
}%footnotesize
\end{Entry}

\clearpage
%%%%%%%%%%%%%%%%%%%%%%%%%%%%%%%%%%%%%%%%%%%%%%%%%%%%%%%%%%%%%%%%%%%%
%
%    Reference: StMcSvtLayerHitCollection
%
%%%%%%%%%%%%%%%%%%%%%%%%%%%%%%%%%%%%%%%%%%%%%%%%%%%%%%%%%%%%%%%%%%%%
\subsection{StMcSvtLayerHitCollection}
\index{SVT hit collection} \index{StMcSvtLayerHitCollection|textbf}
\label{sec:StMcSvtLayerHitCollection}
\textbf{[Renamed]} This class was renamed to \textbf{StMcSvtBarrelHitCollection}
when the SVT hit storage scheme was reorganised from layer--ladder--wafer to
barrel--ladder--wafer order (commit \texttt{f98478bb68}).
See \S\ref{sec:StMcSvtBarrelHitCollection} for the current class reference.
\begin{Entry}
\item[Summary]
    \name{StMcSvtLayerHitCollection} is a container for the
    layers of the SVT.  The actual hits are stored according to
    wafers, 2 levels down the hierarchy.

\item[Synopsis]
    \verb+#include "StMcSvtLayerHitCollection.hh"+\\
    \verb+class StMcSvtLayerHitCollection;+\\

\item[Description]
    \name{StMcSvtLayerHitCollection} is the second step down
    the hierarchy of hits for the SVT.  It provides
    methods to return a particular ladder, and to
    count the number of hits in the layer.

\item[Persistence]
    None

\item[Related Classes]
    \name{StMcSvtLayerHitCollection}
    has a container of type {\tt StMcSvtLadderHitCollection}
    of size 8, each element represents a ladder.  Look in
    \name{StMcSvtHitCollection} \ref{sec:StMcSvtHitCollection},
    \name{StMcSvtLadderHitCollection} \ref{sec:StMcSvtLadderHitCollection},
    and \name{StMcSvtWaferHitCollection} \ref{sec:StMcSvtWaferHitCollection},
    \index{StMcSvtHitCollection}
    \index{StMcSvtLadderHitCollection}
    \index{StMcSvtWaferHitCollection}
    \index{StMcSvtHit}

\item[Public\\ Constructors]
    \verb+StMcSvtLayerHitHitCollection();+\\
    Create an instance of \name{StSvtLayerHitCollection}
    and sets up the internal containers.

\item[Public Member\\ Functions]

    \verb+unsigned long numberOfHits() const;+\\
    Counts the number of hits of the whole SVT.

    \verb+unsigned int  numberOfLadders() const;+\\
    Returns the size of the Array of ladders.  Default is 8.

    \verb+StMcSvtLadderHitCollection*       ladder(unsigned int);+\\
    Returns a pointer to the ladder specified by the argument to the
    function.  Recall that this is for indexing into an array, so the
    argument should go from 0 - (size()-1).  Look in the
    numbering scheme conventions \ref{sec:numberscheme}for more information.

\item[Examples]
{\footnotesize
\begin{verbatim}
//
// Look in StMcSvtWaferHitCollection for an example.
// These classes normally go hand in hand.
\end{verbatim}
}%footnotesize
\end{Entry}

\clearpage
%%%%%%%%%%%%%%%%%%%%%%%%%%%%%%%%%%%%%%%%%%%%%%%%%%%%%%%%%%%%%%%%%%%%
%
%    Reference: StMcSvtWaferHitCollection
%
%%%%%%%%%%%%%%%%%%%%%%%%%%%%%%%%%%%%%%%%%%%%%%%%%%%%%%%%%%%%%%%%%%%%
\subsection{StMcSvtWaferHitCollection}
\index{SVT wafer hit collection} \index{StMcSvtWaferHitCollection|textbf}
\label{sec:StMcSvtWaferHitCollection}
\begin{Entry}
\item[Summary]
    \name{StMcSvtWaferHitCollection} is the class where the SVT hits
    are actually stored.

\item[Synopsis]
    \verb+#include "StMcSvtWaferHitCollection.hh"+\\
    \verb+class StMcSvtWaferHitCollection;+\\

\item[Description]
    \name{StMcSvtWaferHitCollection} holds the SVT hits.

\item[Persistence]
    None

\item[Related Classes]
    \name{StMcSvtWaferHitCollection}
    has a container of type {\tt StSPtrVecMcSvtHit}, this means
    that this is the class that actually owns the hits.  Look
    at the section on containers \ref{sec:containers} for more
    information.    
    \index{StMcSvtHit}
    \index{StSPtrVecMcSvtHit}
\item[Public\\ Constructors]
    \verb+StMcSvtWaferHitCollection();+\\
    Create an instance of \name{StMcSvtWaferHitCollection}
    and sets up the internal container.

\item[Public Member\\ Functions]

    \verb+StSPtrVecMcSvtHit&       hits();+\\
    Returns a reference to the container of the SVT hits.
    
\item[Examples]
{\footnotesize
\begin{verbatim}
//
// Print number of hits in SVT and plot 2D Histogram of x,y of
// Hits
//  
StMcSvtHitCollection* svtHitColl = evt->svtHits();
cout << ``There are :`` << svtHitColl->numberOfHits()
     << ``hits in the SVT'' << endl; 
TH2F* myHist = new TH2F("SVT","pos. of Hits",100, -10, 10, 100, -10, 10)
for (unsigned int ly=0; ly<svtHitColl->numberOfLayers(); ly++) {
  for (unsigned int ld=0; ld<svtHitColl->layer(ly)->numberOfLadders(); ld++) {
    for (unsigned int w=0;
         w<svtHitColl->layer(ly)->ladder(ld)->numberOfWafers();
         w++) {

      StSPtrVecMcSvtHit& hits =
              svtHitColl->layer(ly)->ladder(ld)->wafer(w)->hits();

      for (StMcSvtHitIterator i = hits.begin(); i != hits.end(); i++) {
           currentHit = (*i);
           myHist->Fill(currentHit->position().x(),
                        currentHit->position().y());
      }
    } // wafer  loop
  } // ladder loop 
} // layer  loop
\end{verbatim}
}%footnotesize
\end{Entry}

\clearpage

%%%%%%%%%%%%%%%%%%%%%%%%%%%%%%%%%%%%%%%%%%%%%%%%%%%%%%%%%%%%%%%%%%%%
%
%    Reference: StMcTpcHit
%
%%%%%%%%%%%%%%%%%%%%%%%%%%%%%%%%%%%%%%%%%%%%%%%%%%%%%%%%%%%%%%%%%%%%
\subsection{StMcTpcHit}
\index{TPC hit} \index{StMcTpcHit|textbf}
\label{sec:StMcTpcHit}
\begin{Entry}
\item[Summary]
    \name{StMcTpcHit} represents a monte carlo TPC hit.

\item[Synopsis]
    \verb+#include "StMcTpcHit.hh"+\\
    \verb+class StMcTpcHit;+\\

\item[Description]
    \name{StMcTpcHit} inherits most functionality from \name{StMcHit}.
    For a complete description of the inherited member functions see
    \name{StMcHit} (sec.~\ref{sec:StMcHit}).

    In the \StMcEvent\ data model each track keeps references to the
    associated hits and vice versa. All hits have a pointer to their
    parent track.  This is especially important for the making of the
    associations in \StAssociationMaker.

\item[Persistence]
    None

\item[Related Classes]
    \name{StMcTpcHit} is derived directly from \name{StMcHit}.
    The hits are kept in an instance of \name{StMcTpcHitCollection}
    where they are stored by pointer (see \ref{sec:containers}).
    Each instance of \name{StMcTrack} holds a list of TPC hits
    which belong to that track.
    \index{StMcHit}
    \index{StMcTrack}
    \index{StMcTpcHitCollection}

\item[Public\\ Constructors]
    \verb+StMcTpcHit(const StThreeVectorF& x,const StThreeVectorF& p,+\\
    \verb+      const float de, const float ds, const long key, const long id, StMcTrack* parent);+\\
    
    Create an instance of \name{StMcHit} with position \name{x}, local momentum \name{p},
    energy deposition \name{de}, path length \name{ds}, primary key \name{key},
    volume id \name{id} and parent track \name{parent}. 

\item[Public Member\\ Functions]

    \verb+unsigned long  sector() const;+\\
    Returns the number of the sector [1-24] in which a hit is found. 
    \index{sector, TPC Hit}

    \verb+unsigned long  padrow() const;+\\
    Returns the number of the padrow [1-45] in which a hit is found. 
    \index{padrow, TPC Hit}

\item[Examples]  
{\footnotesize
\begin{verbatim}
//
//  Function which returns the TPC hit
//  that is closest to the
//  start vertex of the track.

StMcTpcHit* getFirstHit(StMcTrack *track)
{
    StThreeVector<float> vtxPosition = track->startVertex();
    StMcTpcHitCollection *hits =
                   track->tpcHits();
    StMcTpcHitIterator iter;
    StMcTpcHit* hit;
    float minDistance = 20*meter;
    for(iter = hits->begin();
        iter != hits->end(); iter++) {
        if (abs(*iter->position() - vtxPosition) < minDistance) {
           minDistance = abs(*iter->position() - vtxPosition);
           hit = *iter;
        }
    }
    return hit;
}

\end{verbatim}
}%footnotesize
\end{Entry}
\clearpage
%%%%%%%%%%%%%%%%%%%%%%%%%%%%%%%%%%%%%%%%%%%%%%%%%%%%%%%%%%%%%%%%%%%%
%
%    Reference: StMcTpcHitCollection
%
%%%%%%%%%%%%%%%%%%%%%%%%%%%%%%%%%%%%%%%%%%%%%%%%%%%%%%%%%%%%%%%%%%%%
\subsection{StMcTpcHitCollection}
\index{TPC hit collection} \index{StMcTpcHitCollection|textbf}
\label{sec:StMcTpcHitCollection}
\begin{Entry}
\item[Summary]
    \name{StMcTpcHitCollection} is a container for the
    sectors of the TPC.  The actual hits are stored according to
    padrows, 2 levels down the hierarchy.

\item[Synopsis]
    \verb+#include "StMcTpcHitCollection.hh"+\\
    \verb+class StMcTpcHitCollection;+\\

\item[Description]
    \name{StMcTpcHitCollection} is the first step down
    the hierarchy of hits for the TPC.  It provides
    methods to return a particular sector, and to
    count the number of hits in all the TPC.

\item[Persistence]
    None

\item[Related Classes]
    \name{StMcTpcHitCollection}
    has a container of type {\tt StMcTpcSectorHitCollection}
    of size 24, each element represents a sector.  Look in
    \name{StMcTpcSectorHitCollection} \ref{sec:StMcTpcSectorHitCollection},
    and \name{StMcTpcPadrowHitCollection} \ref{sec:StMcTpcPadrowHitCollection},
    \index{StMcTpcSectorHitCollection}
    \index{StMcTpcPadrowHitCollection}
    \index{StMcTpcHit}

\item[Public\\ Constructors]
    \verb+StMcTpcHitCollection();+\\
    Create an instance of \name{StTpcHitCollection}
    and sets up the internal containers.

\item[Public Member\\ Functions]

    \verb+bool addHit(StMcTpcHit*);+\\
    Adds a hit to the collection (using \name{StMemoryPool}
    from the StarClassLibrary.

    \verb+unsigned long numberOfHits() const;+\\
    Counts the number of hits of the whole TPC.

    \verb+unsigned int  numberOfSectors() const;+\\
    Returns the size of the Array of sectors.  Default is 24.

    \verb+StMcTpcSectorHitCollection*       sector(unsigned int);+\\
    Returns a pointer to the sector specified by the argument to the
    function.  Recall that this is for indexing into an array, so the
    argument should go from 0 - (size()-1).  Look in the
    numbering scheme conventions \ref{sec:numberscheme}for more information.

\item[Examples]
{\footnotesize
\begin{verbatim}
//
// Look in StMcTpcPadrowHitCollection for an example.
// These classes normally go hand in hand.
\end{verbatim}
}%footnotesize
\end{Entry}

\clearpage
%%%%%%%%%%%%%%%%%%%%%%%%%%%%%%%%%%%%%%%%%%%%%%%%%%%%%%%%%%%%%%%%%%%%
%
%    Reference: StMcTpcPadrowHitCollection
%
%%%%%%%%%%%%%%%%%%%%%%%%%%%%%%%%%%%%%%%%%%%%%%%%%%%%%%%%%%%%%%%%%%%%
\subsection{StMcTpcPadrowHitCollection}
\index{TPC padrow hit collection} \index{StMcTpcPadrowHitCollection|textbf}
\label{sec:StMcTpcPadrowHitCollection}
\begin{Entry}
\item[Summary]
    \name{StMcTpcPadrowHitCollection} is the class where the TPC hits
    are actually stored.

\item[Synopsis]
    \verb+#include "StMcTpcPadrowHitCollection.hh"+\\
    \verb+class StMcTpcPadrowHitCollection;+\\

\item[Description]
    \name{StMcTpcPadrowHitCollection} holds the TPC hits.

\item[Persistence]
    None

\item[Related Classes]
    \name{StMcTpcPadrowHitCollection}
    has a container of type {\tt StSPtrVecMcTpcHit}, this means
    that this is the class that actually owns the hits.  Look
    at the section on containers \ref{sec:containers} for more
    information.    
    \index{StMcTpcHit}
    \index{StSPtrVecMcTpcHit}
\item[Public\\ Constructors]
    \verb+StMcTpcPadrowHitCollection();+\\
    Create an instance of \name{StMcTpcPadrowHitCollection}
    and sets up the internal container.

\item[Public Member\\ Functions]

    \verb+StSPtrVecMcTpcHit&       hits();+\\
    Returns a reference to the container of the TPC hits.
    
\item[Examples]
{\footnotesize
\begin{verbatim}
//
// Print number of hits in TPC and plot 2D Histogram of x,y of
// Hits
//  
StMcTpcHitCollection* tpcHitColl = evt->tpcHits();
cout << ``There are :`` << tpcHitColl->numberOfHits() << ``hits in the TPC'' << endl; 
TH2F* myHist = new TH2F("TPC","x y pos. of Hits",100, -10, 10, 100, -10, 10)

for (unsigned int s=0; s<tpcHitColl->numberOfSectors(); s++) {

  for (unsigned int pr=0; pr<tpcHitColl->sector(s)->numberOfPadrows(); pr++) {

    StSPtrVecMcTpcHit& hits = tpcHitColl->sector(s)->padrow(pr)->hits();

    for (StMcTpcHitIterator i = hits.begin(); i != hits.end(); i++) {

      currentHit = (*i);
      myHist->Fill(currentHit->position().x(), currentHit->position().y());
    }
  } // padrow  loop
} // sector  loop
\end{verbatim}
}%footnotesize
\end{Entry}

\clearpage

%%%%%%%%%%%%%%%%%%%%%%%%%%%%%%%%%%%%%%%%%%%%%%%%%%%%%%%%%%%%%%%%%%%%
%
%    Reference: StMcTpcSectorHitCollection
%
%%%%%%%%%%%%%%%%%%%%%%%%%%%%%%%%%%%%%%%%%%%%%%%%%%%%%%%%%%%%%%%%%%%%
\subsection{StMcTpcSectorHitCollection}
\index{TPC sector hit collection} \index{StMcTpcSectorHitCollection|textbf}
\label{sec:StMcTpcSectorHitCollection}
\begin{Entry}
\item[Summary]
    \name{StMcTpcSectorHitCollection} is a container for the
    sectors of the TPC.  The actual hits are stored according to
    padrows, the next level down the hierarchy.

\item[Synopsis]
    \verb+#include "StMcTpcSectorHitCollection.hh"+\\
    \verb+class StMcTpcSectorHitCollection;+\\

\item[Description]
    \name{StMcTpcSectorHitCollection} is the second step down
    the hierarchy of hits for the TPC.  It provides
    methods to return a particular padrow, and to
    count the number of hits in the sector.

\item[Persistence]
    None

\item[Related Classes]
    \name{StMcTpcSectorHitCollection}
    has a container of type {\tt StMcTpcPadrowHitCollection}
    of size 45, each element represents a padrow.  Look in
    \name{StMcTpcHitCollection} \ref{sec:StMcTpcHitCollection},
    and \name{StMcTpcPadrowHitCollection} \ref{sec:StMcTpcPadrowHitCollection},
    \index{StMcTpcHitCollection}
    \index{StMcTpcPadrowHitCollection}
    \index{StMcTpcHit}

\item[Public\\ Constructors]
    \verb+StMcTpcSectorHitHitCollection();+\\
    Create an instance of \name{StTpcSectorHitCollection}
    and sets up the internal containers.

\item[Public Member\\ Functions]

    \verb+unsigned long numberOfHits() const;+\\
    Counts the number of hits of the whole TPC.

    \verb+unsigned int  numberOfPadrows() const;+\\
    Returns the size of the Array of padrows.  Default is 45.

    \verb+StMcTpcPadrowHitCollection*       padrow(unsigned int);+\\
    Returns a pointer to the padrow specified by the argument to the
    function.  Recall that this is for indexing into an array, so the
    argument should go from 0 - (size()-1).  Look in the
    numbering scheme conventions \ref{sec:numberscheme}for more information.

\item[Examples]
{\footnotesize
\begin{verbatim}
//
// Look in StMcTpcPadrowHitCollection for an example.
// These classes normally go hand in hand.
\end{verbatim}
}%footnotesize
\end{Entry}

\clearpage
%%%%%%%%%%%%%%%%%%%%%%%%%%%%%%%%%%%%%%%%%%%%%%%%%%%%%%%%%%%%%%%%%%%%
%
%    Reference: StMcTrack
%
%%%%%%%%%%%%%%%%%%%%%%%%%%%%%%%%%%%%%%%%%%%%%%%%%%%%%%%%%%%%%%%%%%%%
\subsection{StMcTrack}
\index{StMcTrack|textbf}
\index{monte carlo track}
\label{sec:StMcTrack}
\begin{Entry}
\item[Summary]
    \name{StMcTrack} describes a Monte Carlo track in the STAR detector.

\item[Synopsis]
    \verb+#include "StMcTrack.hh"+\\
    \verb+class StMcTrack;+\\

\item[Description]
    
    \name{StMcTrack} describes a monte carlo track in the STAR
    detector.  
    \name{StMcTrack} provides information on the
    hits which belong to the track, its momentum and the
    vertices of the track.  All
    tracks stored on the g2t\_track table are collected in a
    \name{StSPtrVecMcTrack} which can be obtained through the
    \name{StMcEvent::tracks()} method.  The tracks from the particle
    table are also stored in this container.  For a track that appears
    in both tables, the information is merged so that only one entry
    per track appears in StMcEvent.
    In addition, the information about the particle that produced
    the track is accessible.
    
    The information on the relation between hits and tracks in
    \StMcEvent\ is kept both in \name{StMcTrack} and in \name{StMcHit}.
    Each track knows about the hits it is associated with, and vice versa.
    This was needed in order for the \StAssociationMaker to begin its job by
    building hit associations and from those create the track associations.
    \index{relation btw hits and tracks}
    
\item[Persistence]
    None

\item[Related Classes]
    All monte carlo tracks are kept in an instance of \name{StSPtrVecMcTrack}
    where they are stored by pointer (see sec.~\ref{sec:containers}).
    \index{StMcTrack}
    \index{StSPtrVecMcTrack}

\item[Public\\ Constructors]
    \verb+StMcTrack();+\\
    Constructs an instance of \name{StMcTrack} with all values initialized to 0 (zero).
    Not used.  The tracks are constructed using the g2t\_track table.

\item[Public Member\\ Functions]

    \verb+const StLorentzVectorF&      fourMomentum() const;+\\
    Returns the four-momentum of the track.

    \verb+const StThreeVectorF&        momentum() const;+\\
    Returns the momentum of the track.

    \verb+float                        pt() const;+\\
    Returns pt of the track.

    \verb+float                        rapidity() const;+\\
    Returns rapidity of the track.

    \verb+float                        pseudoRapidity() const;+\\
    Returns pseudoRapidity of the track.

    \verb+StMcVertex* startVertex();+\\
    Returns a pointer to the start vertex of the track.

    \verb+StMcVertex* stopVertex();+\\
    Returns a pointer to the stop vertex of the track.

    \verb+StMcTrack*  parent();+\\
    Returns a pointer to the parent of the track, if it exitsts.  Otherwise, returns a
    null pointer.

    \verb+StPtrVecMcVertex& intermediateVertices();+\\
    Returns a reference to the collection of intermediate vertices of the track.
    Note, that the collection can be empty.

    \verb+StPtrVecMcTpcHit& tpcHits();+\\
    TPC hits belonging to the track.
    The returned container holds the
    hits by pointer. Note, that the collection can be empty.
    This will be the case if the option to load these hits is switched off.

    \verb+StPtrVecMcSvtHit& svtHits();+\\
    SVT hits belonging to the track.
    The returned container holds the
    hits by pointer. Note, that the collection can be empty.
    This will be the case if the option to load these hits is switched off.

    \verb+StPtrVecMcFtpcHit& ftpcHits();+\\
    FTPC hits belonging to the track.
    The returned container holds the
    hits by pointer. Note, that the collection can be empty.
    This will be the case if the option to load these hits is switched off.

    \verb+StPtrVecMcFtpcHit& richHits();+\\
    RICH hits belonging to the track.
    The returned container holds the
    hits by pointer. Note, that the collection can be empty.
    This will be the case if the option to load these hits is switched off.

    \verb+StPtrVecMcRichHit& richHits();+\\
    RICH hits belonging to the track.
    The returned container holds the
    hits by pointer. Note, that the collection can be empty.
    This will be the case if the option to load these hits is switched off.

    \verb+StPtrVecMcCalorimeterHit& bemcHits();+\\
    BEMC hits belonging to the track.
    The returned container holds the
    hits by pointer. Note, that the collection can be empty.
    This will be the case if the option to load these hits is switched off.

    \verb+StPtrVecMcCalorimeterHit& bprsHits();+\\
    Barrel pre-shower hits belonging to the track.
    The returned container holds the
    hits by pointer. Note, that the collection can be empty.
    This will be the case if the option to load these hits is switched off.

    \verb+StPtrVecMcCalorimeterHit& bsmdeHits();+\\
    Barrel shower max-eta hits belonging to the track.
    The returned container holds the
    hits by pointer. Note, that the collection can be empty.
    This will be the case if the option to load these hits is switched off.

    \verb+StPtrVecMcCalorimeterHit& bsmdpHits();+\\
    Barrel shower max-phi hits belonging to the track.
    The returned container holds the
    hits by pointer. Note, that the collection can be empty.
    This will be the case if the option to load these hits is switched off.

    \verb+StParticleDefinition* particleDefinition();+\\
    Returns a pointer to the only instance of the appropriate
    \name{StParticleDefinition} class, which holds
    information on the particle (name, charge, mass, spin, etc).

    \verb+int   isShower();+\\
    Returns 1 if the track is shower, 0 if not.

    \verb+long  geantId();+\\
    GEANT particle ID.

    \verb+long  pdgId();+\\
    PDG code particle ID (when the track comes from the particle table and has
    no GEANT ID).  If the track comes from the g2t table, this will also be set
    but the numerical value will be equal to the geantId value.  (That is what is
    in the g2t table...)

    \verb+long  key();+\\
    Primary key from g2t\_track table.

    \verb+long  eventGenLabel();+\\
    Event Generator Label.  This label should be different from zero if the track
    comes from the event generator.  Useful also in determining the entry in the
    particle table (as index = label - 1).

    \verb+void setFourMomentum(const StLorentzVectorF&);+\\
    Assigns the four-momentum of the track.

    \verb+void setStartVertex(StMcVertex*);+\\
    Assigns the start vertex of the track.

    \verb+void setStopVertex(StMcVertex*);+\\
    Assigns the stop vertex of the track.

    \verb+void setIntermediateVertices(StPtrVecMcVertex&);+\\
    Assigns the intermediate vertex collection of the track.

    \verb+void setTpcHits(StPtrVecMcTpcHit&);+\\
    Assigns the TPC hit collection of the track.

    \verb+void setSvtHits(StPtrVecMcSvtHit&);+\\
    Assigns the SVT hit collection of the track.

    \verb+void setFtpcHits(StPtrVecMcFtpcHit&);+\\
    Assigns the FTPC hit collection of the track.

    \verb+void setRichHits(StPtrVecMcRichHit&);+\\
    Assigns the RICH hit collection of the track.

    \verb+void setBemcHits(StPtrVecMcCalorimeterHit&);+\\
    Assigns the BEMC hit collection of the track.

    \verb+void setBprsHits(StPtrVecMcCalorimeterHit&);+\\
    Assigns the BPRS (pre-shower) hit collection of the track.

    \verb+void setBsmdeHits(StPtrVecMcCalorimeterHit&);+\\
    Assigns the BMDE (shower max-eta) hit collection of the track.

    \verb+void setBsmdpHits(StPtrVecMcCalorimeterHit&);+\\
    Assigns the BMDP (shower max-phi) hit collection of the track.

    \verb+void setShower(char);+\\
    Assigns the shower flag of the track.

    \verb+void setGeantId(long);+\\
    Assigns the GEANT ID of the track.

    \verb+void setPdgId(long);+\\
    Assigns the PDG ID of the track.

    \verb+void setKey(long);+\\
    Assigns the primary key of the track.

    \verb+void setEventGenLabel(long);+\\
    Assigns the event generator label of the track.

    \verb+void setParent(StMcTrack*);+\\
    Assigns the parent track if it exists.

    \verb+void addTpcHit(StMcTpcHit*);+\\
    Adds a TPC hit to the internal hit collection.
    
    \verb+void addFtpcHit(StMcFtpcHit*);+\\
    Adds a FTPC hit to the internal hit collection.

    \verb+void addSvtHit(StMcSvtHit*);+\\
    Adds a SVT hit to the internal hit collection.

    \verb+void addRichHit(StMcRichHit*);+\\
    Adds a RICH hit to the internal hit collection.

    \verb+void addBemcHit(StMcCalorimeterHit*);+\\
    Adds a BEMC hit to the internal hit collection.

    \verb+void addBprsHit(StMcCalorimeterHit*);+\\
    Adds a BPRS hit to the internal hit collection.

    \verb+void addBsmdeHit(StMcCalorimeterHit*);+\\
    Adds a BSMDE hit to the internal hit collection.

    \verb+void addBsmdpHit(StMcCalorimeterHit*);+\\
    Adds a BSMDP hit to the internal hit collection.

    \verb+void removeTpcHit(StMcTpcHit*);+\\
    Removes a TPC hit from the internal hit collection.

    \verb+void removeFtpcHit(StMcFtpcHit*);+\\
    Removes a FTPC hit from the internal hit collection.

    \verb+void removeSvtHit(StMcSvtHit*);+\\
    Removes a SVT hit from the internal hit collection.

    \verb+void removeRichHit(StMcRichHit*);+\\
    Removes a RICH hit from the internal hit collection.

    \verb+void removeCalorimeterHit(StPtrVecMcCalorimeterHit&, StMcCalorimeterHit*);+\\
    Removes a calorimeter hit from the specified vector. Used by the methods below.

    \verb+void removeBemcHit(StMcCalorimeterHit*);+\\
    Removes a Calorimeter hit from the internal BEMC Emc hit collection.  Calls
    removeCalorimeterHit.

    \verb+void removeBprsHit(StMcCalorimeterHit*);+\\
    Removes a Calorimeter hit from the internal BPRS Emc hit collection.  Calls
    removeCalorimeterHit.

    \verb+void removeBsmdeHit(StMcCalorimeterHit*);+\\
    Removes a Calorimeter hit from the internal BSMDE Emc hit collection.  Calls
    removeCalorimeterHit.

    \verb+void removeBsmdpHit(StMcCalorimeterHit*);+\\
    Removes a Calorimeter hit from the internal BSMDP Emc hit collection.  Calls
    removeCalorimeterHit.

\item[Examples]
{\bf Example 1:}
{\footnotesize
\begin{verbatim}
//
//  Calculate the mean pt of all primary
//  tracks.

double meanPtOfTracks(StMcEvent *event)
{
    
    StPtrVecMcTrack& tracks = event->primaryVertex()->daughters();

    StMcTrackIterator iter;
    StMcTrack* theTrack;

    double sumPt = 0;
    int    n = 0;

    for (iter = tracks.begin(); iter != tracks.end(); iter++) {
        theTrack = *iter;       // careful here, pointer collection

        // add up pt
        sumPt += theTrack->pt();
        n++;
    }
    return n ? sumPt/static_cast<double>(n) : 0;
}
\end{verbatim}
}%footnotesize

\end{Entry}
\clearpage
%%%%%%%%%%%%%%%%%%%%%%%%%%%%%%%%%%%%%%%%%%%%%%%%%%%%%%%%%%%%%%%%%%%%
%
%    Reference: StMcVertex
%
%%%%%%%%%%%%%%%%%%%%%%%%%%%%%%%%%%%%%%%%%%%%%%%%%%%%%%%%%%%%%%%%%%%%
\subsection{StMcVertex}
\index{StMcVertex|textbf}
\index{vertex class}
\label{sec:StMcVertex}
\begin{Entry}
\item[Summary]
    \name{StMcVertex} describes a vertex, i.e.~a start or/and end point of one
    or several tracks

\item[Synopsis]
    \verb+#include "StMcVertex.hh"+\\
    \verb+class StMcVertex;+\\

\item[Description]
        \name{StMcVertex} describes a general vertex in the STAR detector.
    It uses the information stored in the \name{g2t\_vertex.idl}
    table.\index{g2t\_vertex.idl} \name{StMcVertex} has a position, a
    GEANT volume, a TOF and a GEANT process. Each vertex also has
    a parent track and 0...n daughter tracks. All are
    referenced by pointers.
    The container of all vertices in the event can be obtained
    through \name{StMcEvent::vertices()}.

\item[Persistence]
    None

\item[Related Classes]
    The StSPtrVecMcVertex container of \name{StMcEvent} owns the
    vertices. This container stores the
    vertices by pointer. 
    \index{tSPtrVecMcVertex}

\item[Public\\ Constructors]
    \verb+StMcVertex();+\\
    Default constructor. Constructs an instance of \name{StMcVertex} with
    values initialized to 0.  Not used.  Construction is done
    using the g2t\_vertex table.

    \verb+StMcVertex(float x, float y, float z);+\\
    Constructs a vertex in position (x, y, z).

\item[Public Member\\ Functions]

    \verb+const StThreeVectorF&       position() const;+\\
    Vertex position in global coordinates.

    \verb+StPtrVecMcTrack& daughters();+\\
    Collection of all daughter tracks. Note that
    the collection can be empty.

    \verb+unsigned int numberOfDaughters();+\\
    Number of daughter tracks. Same as {\tt daughters().size()}.

    \verb+StMcTrack* daughter(unsigned int i);+\\
    Returns the {\tt i}'th daughter track, or a NULL pointer
    if {\tt i} is larger than or equal to the number of stored daughters.
    (Index runs from 0 to {\tt numberOfDaughters()-1}).

    \verb+const StMcTrack* parent();+\\
    Pointer to parent track.
    
    \verb+string geantVolume() const;+\\
    Returns the GEANT volume read from g2t\_vertex table.

    \verb+float tof();+\\
    TOF read from g2t\_vertex table.

    \verb+long geantProcess() const;+\\
    GEANT process read from g2t\_vertex table.

    \verb+long   key() const;+\\
    Primary key, read from g2t\_vertex table.

    \verb+long geantMedium() const;+\\
    GEANT medium read from g2t\_vertex table.

    \verb+long   generatorProcess() const;+\\
    Process from event generator, read from g2t\_vertex table.


    \verb+void setPosition(const StThreeVectorF&);+\\
    Sets vertex position.

    \verb+void setParent(StMcTrack*);+\\
    Sets pointer to parent track.

    \verb+void addDaughter(StMcTrack*); +\\
    Adds a daughter track to the daughter collection.

    \verb+void setGeantVolume(string); +\\
    Sets the GEANT volume.

    \verb+void setTof(float); +\\
    Sets the Time Of Flight.

    \verb+void setGeantProcess(int); +\\
    Sets the GEANT process.

    \verb+void removeDaughter(StMcTrack*);+\\
    Removes a track from the daughter container.

\item[Public Member\\ Operators]
    \verb+int operator==(const StMcVertex&) const;+\\
    Returns true (1) if two instances of \name{StMcVertex} are equal or false (0)
    if otherwise.
    The only data members used for the comparison are the position and the
    GEANT process.

    \verb+int operator!=(const StMcVertex&) const;+\\
    Returns true (1) if two instances of \name{StMcVertex} are not equal or false (0)
    if otherwise.
    This operator is implemented by simply inverting \name{operator==}.

    \verb+int operator!=(const StMcVertex&) const;+\\
    Returns true (1) if two instances of \name{StMcVertex} are not equal or false (0)
    if otherwise.
    This operator is implemented by simply inverting \name{operator==}.

    \verb+ostream&  operator<<(ostream& os, const StMcVertex& v);+\\
    Prints position, geant volume, time of flight and geant process to standard
    output.

\item[Examples] \index{primary vertex}
{\bf Example 1:}
{\footnotesize
\begin{verbatim}
//
//  Function to count # of tracks originating
//  from the primary vertex that have more than
//  10 hits in the TPC.
unsigned long countGoodPrimaryTracks(StMcEvent *event)
{
    unsigned long counter = 0;
    StPtrVecMcTrack& primaryTracks =
                event->primaryVertex()->daughters();
    StMcTrackIterator i;
    StMcTrack* track;
    for (i = primaryTracks.begin(); i != primaryTracks.end(); i++) {
        track = *i;
        if (track->tpcHits()->size() > 10)
            counter++;
    }
    return counter;
}
\end{verbatim}
}%footnotesize

\end{Entry}
\clearpage
%%%%%%%%%%%%%%%%%%%%%%%%%%%%%%%%%%%%%%%%%%%%%%%%%%%%%%%%%%%%%%%%%%%%
%
% The End
%
%%%%%%%%%%%%%%%%%%%%%%%%%%%%%%%%%%%%%%%%%%%%%%%%%%%%%%%%%%%%%%%%%%%%



%% ---- New class references added 2026 ----

%%%%%%%%%%%%%%%%%%%%%%%%%%%%%%%%%%%%%%%%%%%%%%%%%%%%%%%%%%%%%%%%%%%%
%
%    Reference: StMcBTofHit
%
%%%%%%%%%%%%%%%%%%%%%%%%%%%%%%%%%%%%%%%%%%%%%%%%%%%%%%%%%%%%%%%%%%%%
\subsection{StMcBTofHit}
\index{barrel time-of-flight hit} \index{StMcBTofHit|textbf}
\label{sec:StMcBTofHit}
\begin{Entry}
\item[Summary]
    \name{StMcBTofHit} represents a Monte Carlo hit in the Barrel
    Time-of-Flight (BTof) detector.

\item[Synopsis]
    \verb+#include "StMcBTofHit.hh"+\\
    \verb+class StMcBTofHit;+\\

\item[Description]
    \name{StMcBTofHit} stores the tray, module, and cell geometry
    identifiers together with the energy deposited, Geant track path
    length, exact and simulated time-of-flight, collected charge,
    hit position in global coordinates, and a pointer to the parent
    MC track.  Unlike most other STAR MC hit classes it inherits
    directly from \name{StObject} rather than \name{StMcHit}.

\item[Persistence]
    None

\item[Related Classes]
    Hits are owned and accessed through
    \name{StMcBTofHitCollection} (see~\ref{sec:StMcBTofHitCollection}).
    The parent track is of type \name{StMcTrack}.
    \index{StMcBTofHitCollection}
    \index{StMcTrack}

\item[Public\\ Constructors]
    \verb+StMcBTofHit();+\\
    Creates a default-initialised \name{StMcBTofHit}.

    \verb+StMcBTofHit(int tray, int module, int cell,+\\
    \verb+            float dE, float pathLength,+\\
    \verb+            float tof_exact, float tof_measured, float charge);+\\
    Creates an instance with geometry identifiers and physics
    quantities.  The hit position and parent track pointer are
    left at their default values.

    \verb+StMcBTofHit(int tray, int module, int cell,+\\
    \verb+            float dE, float pathLength,+\\
    \verb+            float tof_exact, float tof_measured, float charge,+\\
    \verb+            StThreeVectorF& position, StMcTrack* parentTrack);+\\
    Creates a fully-specified instance including hit position and
    parent track pointer.

\item[Public Member\\ Functions]
    \verb+virtual int tray() const;+\\
    Returns the BTof tray number.\\

    \verb+virtual int module() const;+\\
    Returns the module number within the tray.\\

    \verb+virtual int cell() const;+\\
    Returns the cell number within the module.\\

    \verb+virtual float dE() const;+\\
    Returns the energy deposited in the active material (GeV).\\

    \verb+virtual float pathLength() const;+\\
    Returns the Geant track path length through the active material (cm).\\

    \verb+virtual float time() const;+\\
    Returns the exact Geant time of flight (ns).\\

    \verb+virtual float tof() const;+\\
    Returns the simulated (smeared) time of flight (ns).\\

    \verb+virtual float charge() const;+\\
    Returns the collected charge.\\

    \verb+virtual StThreeVectorF position() const;+\\
    Returns the hit position in global coordinates.\\

    \verb+virtual StMcTrack* parentTrack() const;+\\
    Returns a pointer to the parent MC track.\\

    \verb+virtual int parentTrackId() const;+\\
    Returns the integer identifier of the parent MC track.\\

    \verb+bool sameCell(const StMcBTofHit&) const;+\\
    Returns \texttt{true} if both hits share the same tray, module,
    and cell.\\

    \verb+void operator+=(const StMcBTofHit&);+\\
    Accumulates energy deposit and path length from a hit in the
    same cell into this hit.

\end{Entry}
\clearpage

%%%%%%%%%%%%%%%%%%%%%%%%%%%%%%%%%%%%%%%%%%%%%%%%%%%%%%%%%%%%%%%%%%%%
%
%    Reference: StMcBTofHitCollection
%
%%%%%%%%%%%%%%%%%%%%%%%%%%%%%%%%%%%%%%%%%%%%%%%%%%%%%%%%%%%%%%%%%%%%
\subsection{StMcBTofHitCollection}
\index{detector hit collection} \index{StMcBTofHitCollection|textbf}
\label{sec:StMcBTofHitCollection}
\begin{Entry}
\item[Summary]
    \name{StMcBTofHitCollection} is the top-level container for all
    Barrel Time-of-Flight MC hits in an event.

\item[Synopsis]
    \verb+#include "StMcBTofHitCollection.hh"+\\
    \verb+class StMcBTofHitCollection;+\\

\item[Description]
    \name{StMcBTofHitCollection} owns a flat vector of
    \name{StMcBTofHit} pointers and provides methods to add hits,
    query the hit count, and retrieve the underlying vector.
    The collection stores hits in an \name{StSPtrVecMcBTofHit},
    which manages the lifetimes of the pointed-to objects.

\item[Persistence]
    None

\item[Related Classes]
    Elements of the collection are of type
    \name{StMcBTofHit} (see~\ref{sec:StMcBTofHit}).
    \index{StMcBTofHit}

\item[Public\\ Constructors]
    \verb+StMcBTofHitCollection();+\\
    Creates an empty collection.

\item[Public Member\\ Functions]
    \verb+bool addHit(StMcBTofHit*);+\\
    Adds a hit to the collection; returns \texttt{false} if the
    supplied pointer is null.\\

    \verb+unsigned long numberOfHits() const;+\\
    Returns the number of hits currently in the collection.\\

    \verb+StSPtrVecMcBTofHit& hits();+\\
    Returns a mutable reference to the internal hit vector.\\

    \verb+const StSPtrVecMcBTofHit& hits() const;+\\
    Returns a const reference to the internal hit vector.\\

    \verb+void Clear(const char* opt="");+\\
    Clears and frees all hits owned by the collection.

\end{Entry}
\clearpage

%%%%%%%%%%%%%%%%%%%%%%%%%%%%%%%%%%%%%%%%%%%%%%%%%%%%%%%%%%%%%%%%%%%%
%
%    Reference: StMcMtdHit
%
%%%%%%%%%%%%%%%%%%%%%%%%%%%%%%%%%%%%%%%%%%%%%%%%%%%%%%%%%%%%%%%%%%%%
\subsection{StMcMtdHit}
\index{muon telescope detector hit} \index{StMcMtdHit|textbf}
\label{sec:StMcMtdHit}
\begin{Entry}
\item[Summary]
    \name{StMcMtdHit} represents a Monte Carlo hit in the Muon
    Telescope Detector (MTD).

\item[Synopsis]
    \verb+#include "StMcMtdHit.hh"+\\
    \verb+class StMcMtdHit;+\\

\item[Description]
    \name{StMcMtdHit} extends \name{StMcHit} with the MTD-specific
    geometry identifiers backleg, module, and cell, together with
    the Geant track path length and the parent track identifier.
    It can be constructed either with default values or directly
    from a \texttt{g2t\_mtd\_hit\_st} Geant table row, which is
    the standard path used by \name{StMcEventMaker}.

\item[Persistence]
    None

\item[Related Classes]
    Hits are owned and accessed through
    \name{StMcMtdHitCollection} (see~\ref{sec:StMcMtdHitCollection}).
    \name{StMcMtdHit} inherits position, momentum, energy deposit,
    and time-of-flight from \name{StMcHit}.
    \index{StMcMtdHitCollection}
    \index{StMcHit}

\item[Public\\ Constructors]
    \verb+StMcMtdHit();+\\
    Creates a default-initialised \name{StMcMtdHit}.

    \verb+StMcMtdHit(g2t_mtd_hit_st*);+\\
    Creates an instance of \name{StMcMtdHit} from a Geant table
    row; this is the standard construction path in
    \name{StMcEventMaker}.

\item[Public Member\\ Functions]
    \verb+virtual int backleg() const;+\\
    Returns the MTD backleg number.\\

    \verb+virtual int module() const;+\\
    Returns the module number within the backleg.\\

    \verb+virtual int cell() const;+\\
    Returns the cell number within the module.\\

    \verb+virtual float pathLength() const;+\\
    Returns the Geant track path length through the active material (cm).\\

    \verb+virtual int parentTrackId() const;+\\
    Returns the integer identifier of the parent MC track.\\

    \verb+bool sameCell(const StMcMtdHit&) const;+\\
    Returns \texttt{true} if both hits share the same backleg,
    module, and cell.

\end{Entry}
\clearpage

%%%%%%%%%%%%%%%%%%%%%%%%%%%%%%%%%%%%%%%%%%%%%%%%%%%%%%%%%%%%%%%%%%%%
%
%    Reference: StMcMtdHitCollection
%
%%%%%%%%%%%%%%%%%%%%%%%%%%%%%%%%%%%%%%%%%%%%%%%%%%%%%%%%%%%%%%%%%%%%
\subsection{StMcMtdHitCollection}
\index{detector hit collection} \index{StMcMtdHitCollection|textbf}
\label{sec:StMcMtdHitCollection}
\begin{Entry}
\item[Summary]
    \name{StMcMtdHitCollection} is the top-level container for all
    Muon Telescope Detector MC hits in an event.

\item[Synopsis]
    \verb+#include "StMcMtdHitCollection.hh"+\\
    \verb+class StMcMtdHitCollection;+\\

\item[Description]
    \name{StMcMtdHitCollection} owns a flat vector of
    \name{StMcMtdHit} pointers stored in an
    \name{StSPtrVecMcMtdHit} and provides methods to add hits,
    query the hit count, and retrieve the underlying vector by
    mutable or const reference.

\item[Persistence]
    None

\item[Related Classes]
    Elements of the collection are of type
    \name{StMcMtdHit} (see~\ref{sec:StMcMtdHit}).
    \index{StMcMtdHit}

\item[Public\\ Constructors]
    \verb+StMcMtdHitCollection();+\\
    Creates an empty collection.

\item[Public Member\\ Functions]
    \verb+bool addHit(StMcMtdHit*);+\\
    Adds a hit to the collection; returns \texttt{false} if the
    supplied pointer is null.\\

    \verb+unsigned long numberOfHits() const;+\\
    Returns the number of hits currently in the collection.\\

    \verb+StSPtrVecMcMtdHit& hits();+\\
    Returns a mutable reference to the internal hit vector.\\

    \verb+const StSPtrVecMcMtdHit& hits() const;+\\
    Returns a const reference to the internal hit vector.\\

    \verb+void Clear(const char* opt="");+\\
    Clears and frees all hits owned by the collection.

\end{Entry}
\clearpage

%%%%%%%%%%%%%%%%%%%%%%%%%%%%%%%%%%%%%%%%%%%%%%%%%%%%%%%%%%%%%%%%%%%%
%
%    Reference: StMcTofHit
%
%%%%%%%%%%%%%%%%%%%%%%%%%%%%%%%%%%%%%%%%%%%%%%%%%%%%%%%%%%%%%%%%%%%%
\subsection{StMcTofHit}
\index{time-of-flight hit} \index{StMcTofHit|textbf}
\label{sec:StMcTofHit}
\begin{Entry}
\item[Summary]
    \name{StMcTofHit} represents a Monte Carlo hit for the legacy
    pVPD/TOFr detector (pre-Run~8 era).

\item[Synopsis]
    \verb+#include "StMcTofHit.hh"+\\
    \verb+class StMcTofHit;+\\

\item[Description]
    \name{StMcTofHit} extends \name{StMcHit} with a single
    additional datum: the Geant s-track coordinate along the
    detector strip (\texttt{mStrack}).  All other quantities ---
    position, momentum, energy deposit, path length, and
    time-of-flight --- are inherited from \name{StMcHit}.
    The class can be constructed from explicit quantities or
    directly from a \texttt{g2t\_ctf\_hit\_st} Geant table row,
    which is the standard path used by \name{StMcEventMaker}.

\item[Persistence]
    None

\item[Related Classes]
    Hits are owned and accessed through
    \name{StMcTofHitCollection} (see~\ref{sec:StMcTofHitCollection}).
    \name{StMcTofHit} inherits from \name{StMcHit}.
    \index{StMcTofHitCollection}
    \index{StMcHit}

\item[Public\\ Constructors]
    \verb+StMcTofHit();+\\
    Creates a default-initialised \name{StMcTofHit} with the
    s-track coordinate set to zero.

    \verb+StMcTofHit(const StThreeVectorF& x, const StThreeVectorF& p,+\\
    \verb+           Float_t de=0, Float_t ds=0, Float_t tof=0,+\\
    \verb+           Long_t k=0, Long_t volId=0, StMcTrack* parent=0);+\\
    Creates an instance with explicit hit quantities, forwarding
    all arguments to the \name{StMcHit} base-class constructor and
    initialising the s-track coordinate to zero.

    \verb+StMcTofHit(g2t_ctf_hit_st* pt, Float_t cl_x=0, Float_t cl_t=0);+\\
    Creates an instance of \name{StMcTofHit} from a Geant
    \texttt{g2t\_ctf\_hit\_st} table row; this is the standard
    construction path in \name{StMcEventMaker}.

\item[Public Member\\ Functions]
    \verb+Float_t sTrack() const;+\\
    Returns the Geant s-track coordinate along the detector strip.

\end{Entry}
\clearpage

%%%%%%%%%%%%%%%%%%%%%%%%%%%%%%%%%%%%%%%%%%%%%%%%%%%%%%%%%%%%%%%%%%%%
%
%    Reference: StMcTofHitCollection
%
%%%%%%%%%%%%%%%%%%%%%%%%%%%%%%%%%%%%%%%%%%%%%%%%%%%%%%%%%%%%%%%%%%%%
\subsection{StMcTofHitCollection}
\index{detector hit collection} \index{StMcTofHitCollection|textbf}
\label{sec:StMcTofHitCollection}
\begin{Entry}
\item[Summary]
    \name{StMcTofHitCollection} is the container for all legacy
    pVPD/TOFr MC hits in an event.

\item[Synopsis]
    \verb+#include "StMcTofHitCollection.hh"+\\
    \verb+class StMcTofHitCollection;+\\

\item[Description]
    \name{StMcTofHitCollection} owns a flat vector of
    \name{StMcTofHit} pointers stored in an
    \name{StSPtrVecMcTofHit} and provides methods to add hits,
    query the hit count, and retrieve the underlying vector by
    mutable or const reference.

\item[Persistence]
    None

\item[Related Classes]
    Elements of the collection are of type
    \name{StMcTofHit} (see~\ref{sec:StMcTofHit}).
    \index{StMcTofHit}

\item[Public\\ Constructors]
    \verb+StMcTofHitCollection();+\\
    Creates an empty collection.

\item[Public Member\\ Functions]
    \verb+bool addHit(StMcTofHit*);+\\
    Adds a hit to the collection; returns \texttt{false} if the
    supplied pointer is null.\\

    \verb+unsigned long numberOfHits() const;+\\
    Returns the number of hits currently in the collection.\\

    \verb+StSPtrVecMcTofHit& hits();+\\
    Returns a mutable reference to the internal hit vector.\\

    \verb+const StSPtrVecMcTofHit& hits() const;+\\
    Returns a const reference to the internal hit vector.\\

    \verb+void Clear(const char* opt="");+\\
    Clears and frees all hits owned by the collection.

\end{Entry}
\clearpage

%%%%%%%%%%%%%%%%%%%%%%%%%%%%%%%%%%%%%%%%%%%%%%%%%%%%%%%%%%%%%%%%%%%%
%
%    Reference: StMcCtbHitCollection
%
%%%%%%%%%%%%%%%%%%%%%%%%%%%%%%%%%%%%%%%%%%%%%%%%%%%%%%%%%%%%%%%%%%%%
\subsection{StMcCtbHitCollection}
\index{CTB hit collection} \index{StMcCtbHitCollection|textbf}
\label{sec:StMcCtbHitCollection}
\begin{Entry}
\item[Summary]
    \name{StMcCtbHitCollection} is a flat-array container for all Monte Carlo
    Central Trigger Barrel (CTB) hits in an event.

\item[Synopsis]
    \verb+#include "StMcCtbHitCollection.hh"+\\
    \verb+class StMcCtbHitCollection;+\\

\item[Description]
    \name{StMcCtbHitCollection} holds all \name{StMcCtbHit} objects produced
    in the CTB for a single event.  Unlike the hierarchical collections used
    by detectors such as the TPC, all hits are stored in a single
    \verb+StSPtrVecMcCtbHit+.  The class inherits from \name{StObject} and
    supports ROOT browsing via the \verb+Browse()+ method.

\item[Persistence]
    None

\item[Related Classes]
    Each element of the collection is an \name{StMcCtbHit}
    \ref{sec:StMcCtbHit}.
    \index{StMcCtbHit}

\item[Public\\ Constructors]
    \verb+StMcCtbHitCollection();+\\
    Default constructor; initialises an empty hit vector.

\item[Public Member\\ Functions]
    \verb+bool addHit(StMcCtbHit*);+\\
    Appends the supplied hit pointer to the collection.
    Returns \verb+true+ on success.

    \verb+unsigned long numberOfHits() const;+\\
    Returns the total number of CTB hits stored in the collection.

    \verb+StSPtrVecMcCtbHit& hits();+\\
    \verb+const StSPtrVecMcCtbHit& hits() const;+\\
    Return a reference to (or a const reference to) the internal hit vector.

    \verb+void Clear(const char* opt="");+\\
    Empties the collection, releasing all owned hit objects.

\end{Entry}
\clearpage
%%%%%%%%%%%%%%%%%%%%%%%%%%%%%%%%%%%%%%%%%%%%%%%%%%%%%%%%%%%%%%%%%%%%
%
%    Reference: StMcEtrHit
%
%%%%%%%%%%%%%%%%%%%%%%%%%%%%%%%%%%%%%%%%%%%%%%%%%%%%%%%%%%%%%%%%%%%%
\subsection{StMcEtrHit}
\index{ETR hit} \index{StMcEtrHit|textbf}
\label{sec:StMcEtrHit}
\begin{Entry}
\item[Summary]
    \name{StMcEtrHit} represents a single Monte Carlo hit in the
    End-cap Tagger Ring (ETR).

\item[Synopsis]
    \verb+#include "StMcEtrHit.hh"+\\
    \verb+class StMcEtrHit;+\\

\item[Description]
    \name{StMcEtrHit} inherits from \name{StMcHit} and adds the
    ETR-specific geometric identifiers: the layer number and the
    sector number within that layer.  Instances are normally
    constructed directly from the Geant output table
    \verb+g2t_etr_hit_st+.

\item[Persistence]
    None

\item[Related Classes]
    Hits are owned by \name{StMcEtrHitCollection}
    \ref{sec:StMcEtrHitCollection}.
    The base class \name{StMcHit} supplies position, momentum,
    energy loss, path-length, time-of-flight, key, and parent-track
    information.
    \index{StMcEtrHitCollection}
    \index{StMcHit}

\item[Public\\ Constructors]
    \verb+StMcEtrHit();+\\
    Default constructor; delegates to the \name{StMcHit} default constructor.

    \verb+StMcEtrHit(g2t_etr_hit_st* pt);+\\
    Constructs an \name{StMcEtrHit} from the Geant hit table row pointed to
    by \name{pt}.  This is the standard constructor used by
    \name{StMcEventMaker}.

\item[Public Member\\ Functions]
    \verb+int layer() const;+\\
    Returns the ETR layer number of the hit.

    \verb+int sector() const;+\\
    Returns the sector number within the ETR layer.

    \verb+void Print(Option_t* option="") const;+\\
    Prints a summary of the hit to standard output.  Accessible
    from the ROOT object browser (\verb+*MENU*+).

\end{Entry}
\clearpage
%%%%%%%%%%%%%%%%%%%%%%%%%%%%%%%%%%%%%%%%%%%%%%%%%%%%%%%%%%%%%%%%%%%%
%
%    Reference: StMcEtrHitCollection
%
%%%%%%%%%%%%%%%%%%%%%%%%%%%%%%%%%%%%%%%%%%%%%%%%%%%%%%%%%%%%%%%%%%%%
\subsection{StMcEtrHitCollection}
\index{ETR hit collection} \index{StMcEtrHitCollection|textbf}
\label{sec:StMcEtrHitCollection}
\begin{Entry}
\item[Summary]
    \name{StMcEtrHitCollection} is a container for all Monte Carlo
    End-cap Tagger Ring (ETR) hits in an event.

\item[Synopsis]
    \verb+#include "StMcEtrHitCollection.hh"+\\
    \verb+class StMcEtrHitCollection;+\\

\item[Description]
    \name{StMcEtrHitCollection} holds all \name{StMcEtrHit} objects produced
    in the ETR for a single event, stored in a flat
    \verb+StSPtrVecMcEtrHit+.  The class also exposes the geometry
    constants of the ETR: three layers and twelve sectors per layer.

\item[Persistence]
    None

\item[Related Classes]
    Each element is an \name{StMcEtrHit} \ref{sec:StMcEtrHit}.
    \index{StMcEtrHit}

\item[Public\\ Constructors]
    \verb+StMcEtrHitCollection();+\\
    Default constructor; initialises an empty hit vector.

\item[Enumerations]
    \verb+enum { mNumberOfLayers  = 3  };+\\
    \verb+enum { mNumberOfSectors = 12 };+\\
    Compile-time constants encoding the ETR geometry.

\item[Public Member\\ Functions]
    \verb+bool addHit(StMcEtrHit*);+\\
    Appends the supplied hit pointer to the collection.
    Returns \verb+true+ on success.

    \verb+unsigned long numberOfHits() const;+\\
    Returns the total number of ETR hits stored in the collection.

    \verb+unsigned int numberOfLayers() const;+\\
    Returns the number of ETR layers (currently 3).

    \verb+unsigned int numberOfSectors() const;+\\
    Returns the number of sectors per layer (currently 12).

    \verb+StSPtrVecMcEtrHit& hits();+\\
    \verb+const StSPtrVecMcEtrHit& hits() const;+\\
    Return a reference to (or a const reference to) the internal hit vector.

\end{Entry}
\clearpage
%%%%%%%%%%%%%%%%%%%%%%%%%%%%%%%%%%%%%%%%%%%%%%%%%%%%%%%%%%%%%%%%%%%%
%
%    Reference: StMcFgtHit
%
%%%%%%%%%%%%%%%%%%%%%%%%%%%%%%%%%%%%%%%%%%%%%%%%%%%%%%%%%%%%%%%%%%%%
\subsection{StMcFgtHit}
\index{FGT hit} \index{StMcFgtHit|textbf}
\label{sec:StMcFgtHit}
\begin{Entry}
\item[Summary]
    \name{StMcFgtHit} represents a single Monte Carlo hit in the
    Forward GEM Tracker (FGT).

\item[Synopsis]
    \verb+#include "StMcFgtHit.hh"+\\
    \verb+class StMcFgtHit;+\\

\item[Description]
    \name{StMcFgtHit} inherits from \name{StMcHit} and adds the
    FGT-specific identifiers: the layer (disk) number and the
    quadrant.  The layer and quadrant are encoded in the
    \verb+volume_id+ field of the Geant hit table and are decoded
    by the accessor methods.

\item[Persistence]
    None

\item[Related Classes]
    Hits are owned by \name{StMcFgtLayerHitCollection}
    \ref{sec:StMcFgtLayerHitCollection}, which is in turn owned by
    \name{StMcFgtHitCollection} \ref{sec:StMcFgtHitCollection}.
    The base class \name{StMcHit} provides position, momentum,
    energy loss, path-length, time-of-flight, key, and parent-track
    accessors.
    \index{StMcFgtHitCollection}
    \index{StMcFgtLayerHitCollection}
    \index{StMcHit}

\item[Public\\ Constructors]
    \verb+StMcFgtHit();+\\
    Default constructor; delegates to the \name{StMcHit} default constructor.

    \verb+StMcFgtHit(const StThreeVectorF& x, const StThreeVectorF& p,+\\
    \verb+           Float_t de, Float_t ds, Float_t tof, Long_t k,+\\
    \verb+           Long_t volId, StMcTrack* parent);+\\
    Constructs an \name{StMcFgtHit} from explicit hit parameters.

    \verb+StMcFgtHit(g2t_fgt_hit_st* pt);+\\
    Constructs an \name{StMcFgtHit} from the Geant FGT hit table row
    pointed to by \name{pt}.  This is the standard constructor used
    by \name{StMcEventMaker}.

\item[Public Member\\ Functions]
    \verb+ULong_t layer() const;+\\
    Returns the FGT layer (disk) number decoded from the volume identifier.

    \verb+ULong_t quad() const;+\\
    Returns the FGT quadrant number decoded from the volume identifier.

    \verb+virtual void Print(Option_t* option="") const;+\\
    Prints a summary of the hit to standard output.  Accessible
    from the ROOT object browser (\verb+*MENU*+).

\end{Entry}
\clearpage
%%%%%%%%%%%%%%%%%%%%%%%%%%%%%%%%%%%%%%%%%%%%%%%%%%%%%%%%%%%%%%%%%%%%
%
%    Reference: StMcFgtHitCollection
%
%%%%%%%%%%%%%%%%%%%%%%%%%%%%%%%%%%%%%%%%%%%%%%%%%%%%%%%%%%%%%%%%%%%%
\subsection{StMcFgtHitCollection}
\index{FGT hit collection} \index{StMcFgtHitCollection|textbf}
\label{sec:StMcFgtHitCollection}
\begin{Entry}
\item[Summary]
    \name{StMcFgtHitCollection} is the top-level container for all Monte
    Carlo Forward GEM Tracker (FGT) hits in an event, organised by layer.

\item[Synopsis]
    \verb+#include "StMcFgtHitCollection.hh"+\\
    \verb+class StMcFgtHitCollection;+\\

\item[Description]
    \name{StMcFgtHitCollection} owns an array of
    \name{StMcFgtLayerHitCollection} objects, one per FGT disk (layer).
    There are nine layers in total, corresponding to the nine FGT disk
    positions as defined in \verb+StFgtGeom+.  Hits are routed to the
    correct layer sub-collection by \verb+addHit()+.  The class inherits
    from \name{StObject}.

\item[Persistence]
    None

\item[Related Classes]
    Each layer sub-collection is an \name{StMcFgtLayerHitCollection}
    \ref{sec:StMcFgtLayerHitCollection}.
    Each hit within a layer is an \name{StMcFgtHit} \ref{sec:StMcFgtHit}.
    \index{StMcFgtLayerHitCollection}
    \index{StMcFgtHit}

\item[Public\\ Constructors]
    \verb+StMcFgtHitCollection();+\\
    Default constructor; initialises the nine layer sub-collections.

\item[Enumerations]
    \verb+enum { mNumberOfLayers = 9 };+\\
    Compile-time constant for the number of FGT layers (disks).

\item[Public Member\\ Functions]
    \verb+bool addHit(StMcFgtHit*);+\\
    Routes the supplied hit to the appropriate layer sub-collection.
    Returns \verb+true+ on success.

    \verb+unsigned long numberOfHits() const;+\\
    Returns the total number of FGT hits summed across all layers.

    \verb+unsigned int numberOfLayers() const;+\\
    Returns the number of FGT layers (currently 9).

    \verb+StMcFgtLayerHitCollection* layer(unsigned int);+\\
    \verb+const StMcFgtLayerHitCollection* layer(unsigned int) const;+\\
    Return a pointer to (or a const pointer to) the
    \name{StMcFgtLayerHitCollection} for the specified layer index
    (0-based).

\end{Entry}
\clearpage
%%%%%%%%%%%%%%%%%%%%%%%%%%%%%%%%%%%%%%%%%%%%%%%%%%%%%%%%%%%%%%%%%%%%
%
%    Reference: StMcFgtLayerHitCollection
%
%%%%%%%%%%%%%%%%%%%%%%%%%%%%%%%%%%%%%%%%%%%%%%%%%%%%%%%%%%%%%%%%%%%%
\subsection{StMcFgtLayerHitCollection}
\index{FGT layer hit collection} \index{StMcFgtLayerHitCollection|textbf}
\label{sec:StMcFgtLayerHitCollection}
\begin{Entry}
\item[Summary]
    \name{StMcFgtLayerHitCollection} is the per-layer sub-collection of
    Monte Carlo FGT hits, holding all hits recorded in a single FGT disk.

\item[Synopsis]
    \verb+#include "StMcFgtLayerHitCollection.hh"+\\
    \verb+class StMcFgtLayerHitCollection;+\\

\item[Description]
    \name{StMcFgtLayerHitCollection} stores \name{StMcFgtHit} objects
    belonging to one FGT disk layer in a
    \verb+StSPtrVecMcFgtHit+.  Nine instances of this class are owned
    by \name{StMcFgtHitCollection}, one per layer.  The class inherits
    from \name{StObject} and supports ROOT browsing.

\item[Persistence]
    None

\item[Related Classes]
    Nine instances are owned by \name{StMcFgtHitCollection}
    \ref{sec:StMcFgtHitCollection}.
    Each element is an \name{StMcFgtHit} \ref{sec:StMcFgtHit}.
    \index{StMcFgtHitCollection}
    \index{StMcFgtHit}

\item[Public\\ Constructors]
    \verb+StMcFgtLayerHitCollection();+\\
    Default constructor; initialises an empty hit vector.

\item[Public Member\\ Functions]
    \verb+unsigned long numberOfHits() const;+\\
    Returns the number of hits in this layer.

    \verb+StSPtrVecMcFgtHit& hits();+\\
    \verb+const StSPtrVecMcFgtHit& hits() const;+\\
    Return a reference to (or a const reference to) the internal hit vector.

    \verb+void Clear(const char* opt="");+\\
    Empties the layer collection, releasing all owned hit objects.

\end{Entry}
\clearpage
%%%%%%%%%%%%%%%%%%%%%%%%%%%%%%%%%%%%%%%%%%%%%%%%%%%%%%%%%%%%%%%%%%%%
%
%    Reference: StMcHitT
%
%%%%%%%%%%%%%%%%%%%%%%%%%%%%%%%%%%%%%%%%%%%%%%%%%%%%%%%%%%%%%%%%%%%%
\subsection{StMcHitT\textless{}T\textgreater{}}
\index{StMcHitT|textbf}
\label{sec:StMcHitT}
\begin{Entry}
\item[Summary]
    \name{StMcHitT\textless{}T\textgreater{}} is a template adapter that wraps a
    single row of a Geant hit table (a \verb+g2t_*_hit_st+ struct) and
    exposes it through the \name{StMcHitI} interface.

\item[Synopsis]
    \verb+#include "StMcHitT.hh"+\\
    \verb+template <class T> class StMcHitT;+\\

\item[Description]
    \name{StMcHitT\textless{}T\textgreater{}} inherits from both
    \verb+StEventObject<T>+ (which holds a pointer or value of type
    \verb+T+) and from the pure interface \name{StMcHitI}.
    It provides concrete implementations of all \name{StMcHitI} accessors
    by reading directly from the underlying table row.
    This design avoids duplicating the accessor implementations across every
    concrete MC hit class that maps one-to-one onto a Geant table.

    A full explicit specialisation for \verb+g2t_emc_hit_st*+ is provided
    in the same header to handle the reduced field set of the EMC hit table.

    Concrete hit classes (e.g.\ \name{StMcFgtHit}, \name{StMcEtrHit})
    inherit from \name{StMcHit} rather than from \name{StMcHitT} directly;
    \name{StMcHitT} is used where a hit class maps directly and completely
    onto a single Geant table row.

\item[Persistence]
    None

\item[Related Classes]
    \name{StMcHitI} defines the pure virtual interface satisfied by this
    template.
    \verb+StEventObject<T>+ manages the underlying table-row storage.
    \index{StMcHitI}

\item[Template\\ Parameter]
    \verb+T+ --- a pointer to a Geant hit table row struct
    (e.g.\ \verb+g2t_fgt_hit_st*+).  The struct must expose at minimum
    the fields \verb+x[3]+, \verb+p[3]+, \verb+de+, \verb+ds+,
    \verb+tof+, \verb+id+, \verb+volume_id+, and \verb+track_p+.

\item[Public\\ Constructors]
    \verb+StMcHitT(T data);+\\
    Constructs the adapter from an existing table-row pointer \name{data},
    forwarding it to \verb+StEventObject<T>+.

\item[Public Member\\ Functions]
    \verb+virtual float x()       const;+\\
    \verb+virtual float y()       const;+\\
    \verb+virtual float z()       const;+\\
    Return the $x$, $y$, and $z$ hit position coordinates (cm) read from
    \verb+fData->x[0..2]+.

    \verb+virtual float px()      const;+\\
    \verb+virtual float py()      const;+\\
    \verb+virtual float pz()      const;+\\
    Return the $x$, $y$, and $z$ local momentum components (GeV/$c$) read
    from \verb+fData->p[0..2]+.

    \verb+virtual float dE()      const;+\\
    Returns the energy deposited in the hit volume (GeV), from
    \verb+fData->de+.

    \verb+virtual float dS()      const;+\\
    Returns the path length through the hit volume (cm), from
    \verb+fData->ds+.

    \verb+virtual float tof()     const;+\\
    Returns the time of flight to the hit (ns), from \verb+fData->tof+.

    \verb+virtual long  key()     const;+\\
    Returns the Geant hit identifier, from \verb+fData->id+.

    \verb+virtual long  volumeId() const;+\\
    Returns the encoded volume identifier, from \verb+fData->volume_id+.

    \verb+virtual long  parentTrackIndex() const;+\\
    Returns the index of the parent Geant track, from
    \verb+fData->track_p+.

    \verb+int operator==(const StMcHitT&) const;+\\
    \verb+int operator!=(const StMcHitT&) const;+\\
    Equality and inequality operators, delegating to
    \verb+StEventObject<T>::operator==()+.

\end{Entry}

%%%%%%%%%%%%%%%%%%%%%%%%%%%%%%%%%%%%%%%%%%%%%%%%%%%%%%%%%%%%%%%%%%%%
%
%    Reference: StMcIstHit
%
%%%%%%%%%%%%%%%%%%%%%%%%%%%%%%%%%%%%%%%%%%%%%%%%%%%%%%%%%%%%%%%%%%%%
\subsection{StMcIstHit}
\index{IST hit} \index{StMcIstHit|textbf}
\label{sec:StMcIstHit}
\begin{Entry}
\item[Summary]
    \name{StMcIstHit} represents a monte carlo IST (Inner Silicon Tracker) hit.

\item[Synopsis]
    \verb+#include "StMcIstHit.hh"+\\
    \verb+class StMcIstHit;+\\

\item[Description]
    \name{StMcIstHit} inherits most functionality from \name{StMcHit}.
    For a complete description of the inherited member functions see
    \name{StMcHit} (sec.~\ref{sec:StMcHit}).

    In the \StMcEvent\ data model each track keeps references to the
    associated hits and vice versa.  All hits have a pointer to their
    parent track.  This is especially important for the making of the
    associations in \StAssociationMaker.

\item[Persistence]
    None

\item[Related Classes]
    \name{StMcIstHit} is derived directly from \name{StMcHit}.
    The way to retrieve hits for a given ladder and sensor starting
    from the top-level \name{StMcEvent} pointer is:\\
    \verb+mcEvt->istHitCollection()->layer(ly)->hits()+\\
    (using the layer-based collection), or via the ladder/sensor hierarchy
    introduced in 2015.
    The hits are kept in an instance of \name{StSPtrVecMcIstHit}
    where they are stored by pointer (see \ref{sec:containers}).
    Each instance of \name{StMcTrack} holds a list of IST hits
    which belong to that track.
    \index{StMcHit}
    \index{StMcTrack}
    \index{StMcIstHitCollection}

\item[Public\\ Constructors]
    \verb+StMcIstHit();+\\
    Default constructor.

    \verb+StMcIstHit(const StThreeVectorF& x, const StThreeVectorF& p,+\\
    \verb+           Float_t de, Float_t ds, Float_t tof,+\\
    \verb+           Long_t k, Long_t volId, StMcTrack* parent);+\\
    Create an instance of \name{StMcIstHit} with position \name{x},
    local momentum \name{p}, energy deposition \name{de}, path length
    \name{ds}, time of flight \name{tof}, primary key \name{k},
    volume id \name{volId} and parent track \name{parent}.

    \verb+StMcIstHit(g2t_ist_hit_st* pt);+\\
    Construct from a GEANT \texttt{g2t\_ist\_hit} table row.

\item[Public Member\\ Functions]

    \verb+ULong_t layer() const;+\\
    Returns the layer number; always returns 1 for the IST.
    \index{layer, IST}

    \verb+ULong_t ladder() const;+\\
    Returns the ladder number decoded from the volume id.
    \index{ladder, IST}

    \verb+ULong_t wafer() const;+\\
    Returns the wafer (sensor) number decoded from the volume id.
    \index{wafer, IST}

    \verb+ULong_t side();+\\
    Returns the side of the ladder (1\,=\,inner; 2\,=\,outer) decoded
    from the volume id.
    \index{side, IST}

\item[Examples]
{\footnotesize
\begin{verbatim}
//
//  Print IST hits belonging to a track.
//
void printIstHits(StMcTrack* track)
{
    StPtrVecMcIstHit& hits = track->istHits();
    for (StMcIstHitIterator it = hits.begin();
         it != hits.end(); ++it)
        cout << *(*it) << endl;
}
\end{verbatim}
}%footnotesize
\end{Entry}
\clearpage

%%%%%%%%%%%%%%%%%%%%%%%%%%%%%%%%%%%%%%%%%%%%%%%%%%%%%%%%%%%%%%%%%%%%
%
%    Reference: StMcIstHitCollection
%
%%%%%%%%%%%%%%%%%%%%%%%%%%%%%%%%%%%%%%%%%%%%%%%%%%%%%%%%%%%%%%%%%%%%
\subsection{StMcIstHitCollection}
\index{IST hit collection} \index{StMcIstHitCollection|textbf}
\label{sec:StMcIstHitCollection}
\begin{Entry}
\item[Summary]
    \name{StMcIstHitCollection} is the top-level container for IST monte
    carlo hits, organised by layer.

\item[Synopsis]
    \verb+#include "StMcIstHitCollection.hh"+\\
    \verb+class StMcIstHitCollection;+\\

\item[Description]
    \name{StMcIstHitCollection} is the first step down the hierarchy of hits
    for the IST.  It provides methods to return a particular layer
    sub-collection and to count the total number of hits across all IST layers.
    The collection holds an array of 24 \name{StMcIstLayerHitCollection}
    objects.

\item[Persistence]
    None

\item[Related Classes]
    \name{StMcIstHitCollection} contains an array of
    {\tt StMcIstLayerHitCollection} objects (size~24).  See
    \name{StMcIstLayerHitCollection} \ref{sec:StMcIstLayerHitCollection}
    and \name{StMcIstHit} \ref{sec:StMcIstHit}.
    \index{StMcIstLayerHitCollection}
    \index{StMcIstHit}

\item[Public\\ Constructors]
    \verb+StMcIstHitCollection();+\\
    Create an instance of \name{StMcIstHitCollection} and
    initialise the internal layer sub-collections.

\item[Public Member\\ Functions]

    \verb+bool addHit(StMcIstHit*);+\\
    Add a hit to the appropriate layer sub-collection.

    \verb+unsigned long numberOfHits() const;+\\
    Count the total number of hits across all layers.

    \verb+unsigned int  numberOfLayers() const;+\\
    Return the number of layer sub-collections (24).

    \verb+StMcIstLayerHitCollection*       layer(unsigned int);+\\
    \verb+const StMcIstLayerHitCollection* layer(unsigned int) const;+\\
    Return the mutable or const layer sub-collection at 0-based index
    \name{i}.

\item[Examples]
{\footnotesize
\begin{verbatim}
//
// Count total IST hits and print per-layer totals.
//
StMcIstHitCollection* coll = mcEvt->istHitCollection();
cout << "Total IST hits: " << coll->numberOfHits() << endl;
for (unsigned int ly = 0; ly < coll->numberOfLayers(); ++ly)
    cout << "  layer " << ly << ": "
         << coll->layer(ly)->numberOfHits() << endl;
\end{verbatim}
}%footnotesize
\end{Entry}
\clearpage

%%%%%%%%%%%%%%%%%%%%%%%%%%%%%%%%%%%%%%%%%%%%%%%%%%%%%%%%%%%%%%%%%%%%
%
%    Reference: StMcIstLadderHitCollection
%
%%%%%%%%%%%%%%%%%%%%%%%%%%%%%%%%%%%%%%%%%%%%%%%%%%%%%%%%%%%%%%%%%%%%
\subsection{StMcIstLadderHitCollection}
\index{IST ladder hit collection} \index{StMcIstLadderHitCollection|textbf}
\label{sec:StMcIstLadderHitCollection}
\begin{Entry}
\item[Summary]
    \name{StMcIstLadderHitCollection} is a container for IST hits in one
    ladder, organised by sensor.

\item[Synopsis]
    \verb+#include "StMcIstLadderHitCollection.hh"+\\
    \verb+class StMcIstLadderHitCollection;+\\

\item[Description]
    \name{StMcIstLadderHitCollection} forms the second level of the
    ladder--sensor hit hierarchy introduced for the IST in 2015.  Each
    ladder holds an array of up to 6
    \name{StMcIstSensorHitCollection} objects.

\item[Persistence]
    None

\item[Related Classes]
    \name{StMcIstLadderHitCollection} contains an array of
    {\tt StMcIstSensorHitCollection} objects (size~6).  See
    \name{StMcIstSensorHitCollection} \ref{sec:StMcIstSensorHitCollection}
    and \name{StMcIstHit} \ref{sec:StMcIstHit}.
    \index{StMcIstSensorHitCollection}
    \index{StMcIstHit}

\item[Public\\ Constructors]
    \verb+StMcIstLadderHitCollection();+\\
    Create an instance of \name{StMcIstLadderHitCollection} and
    initialise the internal sensor sub-collections.

\item[Public Member\\ Functions]

    \verb+unsigned long numberOfHits() const;+\\
    Count the total number of hits across all sensors in this ladder.

    \verb+unsigned int  numberOfSensors() const;+\\
    Return the number of sensor sub-collections (6).

    \verb+StMcIstSensorHitCollection*       sensor(unsigned int);+\\
    \verb+const StMcIstSensorHitCollection* sensor(unsigned int) const;+\\
    Return the mutable or const sensor sub-collection at 0-based index
    \name{i}.

\item[Examples]
{\footnotesize
\begin{verbatim}
//
// Look in StMcIstSensorHitCollection for an example.
// These classes normally go hand in hand.
\end{verbatim}
}%footnotesize
\end{Entry}
\clearpage

%%%%%%%%%%%%%%%%%%%%%%%%%%%%%%%%%%%%%%%%%%%%%%%%%%%%%%%%%%%%%%%%%%%%
%
%    Reference: StMcIstLayerHitCollection
%
%%%%%%%%%%%%%%%%%%%%%%%%%%%%%%%%%%%%%%%%%%%%%%%%%%%%%%%%%%%%%%%%%%%%
\subsection{StMcIstLayerHitCollection}
\index{IST layer hit collection} \index{StMcIstLayerHitCollection|textbf}
\label{sec:StMcIstLayerHitCollection}
\begin{Entry}
\item[Summary]
    \name{StMcIstLayerHitCollection} is the leaf node in the
    layer-based IST hit hierarchy; it directly stores the IST hits for
    one layer.

\item[Synopsis]
    \verb+#include "StMcIstLayerHitCollection.hh"+\\
    \verb+class StMcIstLayerHitCollection;+\\

\item[Description]
    \name{StMcIstLayerHitCollection} is the second (and final) step down
    the layer-based hierarchy for IST hits.  It stores the actual
    \name{StMcIstHit} pointers in a \name{StSPtrVecMcIstHit} container.

\item[Persistence]
    None

\item[Related Classes]
    \name{StMcIstLayerHitCollection} contains a
    {\tt StSPtrVecMcIstHit}, which owns the hit objects.  See the
    section on containers \ref{sec:containers} for more information.
    The parent collection is
    \name{StMcIstHitCollection} \ref{sec:StMcIstHitCollection}.
    \index{StMcIstHitCollection}
    \index{StMcIstHit}
    \index{StSPtrVecMcIstHit}

\item[Public\\ Constructors]
    \verb+StMcIstLayerHitCollection();+\\
    Create an instance of \name{StMcIstLayerHitCollection} and
    initialise the internal hit container.

\item[Public Member\\ Functions]

    \verb+unsigned long numberOfHits() const;+\\
    Return the number of hits stored in this layer.

    \verb+StSPtrVecMcIstHit&       hits();+\\
    \verb+const StSPtrVecMcIstHit& hits() const;+\\
    Return a reference to the container of IST hits for this layer.

\item[Examples]
{\footnotesize
\begin{verbatim}
//
// Iterate over all hits in a given IST layer.
//
StMcIstHitCollection* coll = mcEvt->istHitCollection();
StSPtrVecMcIstHit& hits = coll->layer(0)->hits();
for (StMcIstHitIterator it = hits.begin();
     it != hits.end(); ++it)
    cout << *(*it) << endl;
\end{verbatim}
}%footnotesize
\end{Entry}
\clearpage

%%%%%%%%%%%%%%%%%%%%%%%%%%%%%%%%%%%%%%%%%%%%%%%%%%%%%%%%%%%%%%%%%%%%
%
%    Reference: StMcIstSensorHitCollection
%
%%%%%%%%%%%%%%%%%%%%%%%%%%%%%%%%%%%%%%%%%%%%%%%%%%%%%%%%%%%%%%%%%%%%
\subsection{StMcIstSensorHitCollection}
\index{IST sensor hit collection} \index{StMcIstSensorHitCollection|textbf}
\label{sec:StMcIstSensorHitCollection}
\begin{Entry}
\item[Summary]
    \name{StMcIstSensorHitCollection} is the leaf node in the
    ladder--sensor IST hit hierarchy; it directly stores the IST hits for
    one sensor.

\item[Synopsis]
    \verb+#include "StMcIstSensorHitCollection.hh"+\\
    \verb+class StMcIstSensorHitCollection;+\\

\item[Description]
    \name{StMcIstSensorHitCollection} stores the actual
    \name{StMcIstHit} pointers in a \name{StSPtrVecMcIstHit} container
    for a single IST sensor.  It is part of the ladder--sensor hierarchy
    introduced in 2015.

\item[Persistence]
    None

\item[Related Classes]
    \name{StMcIstSensorHitCollection} contains a
    {\tt StSPtrVecMcIstHit}, which owns the hit objects.  See the
    section on containers \ref{sec:containers}.
    The parent collection is
    \name{StMcIstLadderHitCollection} \ref{sec:StMcIstLadderHitCollection}.
    \index{StMcIstLadderHitCollection}
    \index{StMcIstHit}
    \index{StSPtrVecMcIstHit}

\item[Public\\ Constructors]
    \verb+StMcIstSensorHitCollection();+\\
    Create an instance of \name{StMcIstSensorHitCollection} and
    initialise the internal hit container.

\item[Public Member\\ Functions]

    \verb+StSPtrVecMcIstHit&       hits();+\\
    \verb+const StSPtrVecMcIstHit& hits() const;+\\
    Return a reference to the container of IST hits for this sensor.

\item[Examples]
{\footnotesize
\begin{verbatim}
//
// Look in StMcIstLadderHitCollection for an example.
// These classes normally go hand in hand.
\end{verbatim}
}%footnotesize
\end{Entry}
\clearpage

%%%%%%%%%%%%%%%%%%%%%%%%%%%%%%%%%%%%%%%%%%%%%%%%%%%%%%%%%%%%%%%%%%%%
%
%    Reference: StMcPxlHit
%
%%%%%%%%%%%%%%%%%%%%%%%%%%%%%%%%%%%%%%%%%%%%%%%%%%%%%%%%%%%%%%%%%%%%
\subsection{StMcPxlHit}
\index{PXL hit} \index{StMcPxlHit|textbf}
\label{sec:StMcPxlHit}
\begin{Entry}
\item[Summary]
    \name{StMcPxlHit} represents a monte carlo PXL (Pixel) detector hit.

\item[Synopsis]
    \verb+#include "StMcPxlHit.hh"+\\
    \verb+class StMcPxlHit;+\\

\item[Description]
    \name{StMcPxlHit} inherits most functionality from \name{StMcHit}.
    For a complete description of the inherited member functions see
    \name{StMcHit} (sec.~\ref{sec:StMcHit}).

    In the \StMcEvent\ data model each track keeps references to the
    associated hits and vice versa.  All hits have a pointer to their
    parent track.  This is especially important for the making of the
    associations in \StAssociationMaker.

    The PXL detector is organised into sectors (10), ladders (4 per
    sector), and sensors (10 per ladder).  The geometry numbering
    scheme is described at
    \texttt{drupal.star.bnl.gov/STAR/system/files/HFT\%20numbering\%20scheme\_v5\_0.pdf}.

\item[Persistence]
    None

\item[Related Classes]
    \name{StMcPxlHit} is derived directly from \name{StMcHit}.
    The way to retrieve hits for a given sector, ladder, and sensor
    starting from the top-level \name{StMcEvent} pointer is:\\
    \verb+mcEvt->pxlHitCollection()->sector(s)->ladder(ld)->sensor(sn)->hits()+\\
    Each instance of \name{StMcTrack} holds a list of PXL hits
    which belong to that track.
    \index{StMcHit}
    \index{StMcTrack}
    \index{StMcPxlHitCollection}

\item[Public\\ Constructors]
    \verb+StMcPxlHit();+\\
    Default constructor.

    \verb+StMcPxlHit(const StThreeVectorF& x, const StThreeVectorF& p,+\\
    \verb+           float de, float ds, float tof,+\\
    \verb+           long int k, long int volId, StMcTrack* parent);+\\
    Create an instance of \name{StMcPxlHit} with position \name{x},
    local momentum \name{p}, energy deposition \name{de}, path length
    \name{ds}, time of flight \name{tof}, primary key \name{k},
    volume id \name{volId} and parent track \name{parent}.

    \verb+StMcPxlHit(g2t_pix_hit_st* pt);+\\
    Construct from a GEANT \texttt{g2t\_pix\_hit} table row.

\item[Public Member\\ Functions]

    \verb+int sector() const;+\\
    Returns the sector number (1--10, numbered clockwise when viewed
    from East).
    \index{sector, PXL}

    \verb+int ladder() const;+\\
    Returns the ladder number (1--4, numbered clockwise when viewed from
    East; ladder~4 is the inner ladder).
    \index{ladder, PXL}

    \verb+int sensor() const;+\\
    Returns the sensor number (1--10, numbered from East to West).
    \index{sensor, PXL}

\item[Examples]
{\footnotesize
\begin{verbatim}
//
//  Print PXL hits belonging to a track.
//
void printPxlHits(StMcTrack* track)
{
    StPtrVecMcPxlHit& hits = track->pxlHits();
    for (StMcPxlHitIterator it = hits.begin();
         it != hits.end(); ++it)
        cout << *(*it) << endl;
}
\end{verbatim}
}%footnotesize
\end{Entry}
\clearpage

%%%%%%%%%%%%%%%%%%%%%%%%%%%%%%%%%%%%%%%%%%%%%%%%%%%%%%%%%%%%%%%%%%%%
%
%    Reference: StMcPxlHitCollection
%
%%%%%%%%%%%%%%%%%%%%%%%%%%%%%%%%%%%%%%%%%%%%%%%%%%%%%%%%%%%%%%%%%%%%
\subsection{StMcPxlHitCollection}
\index{PXL hit collection} \index{StMcPxlHitCollection|textbf}
\label{sec:StMcPxlHitCollection}
\begin{Entry}
\item[Summary]
    \name{StMcPxlHitCollection} is the top-level container for PXL monte
    carlo hits, organised by sector.

\item[Synopsis]
    \verb+#include "StMcPxlHitCollection.hh"+\\
    \verb+class StMcPxlHitCollection;+\\

\item[Description]
    \name{StMcPxlHitCollection} is the first step down the hierarchy of hits
    for the PXL detector.  It provides methods to return a particular sector
    sub-collection and to count the total number of hits across all sectors.
    The collection holds an array of 10 \name{StMcPxlSectorHitCollection}
    objects, one per sector.

\item[Persistence]
    None

\item[Related Classes]
    \name{StMcPxlHitCollection} contains an array of
    {\tt StMcPxlSectorHitCollection} objects (size~10).  See
    \name{StMcPxlSectorHitCollection} \ref{sec:StMcPxlSectorHitCollection},
    \name{StMcPxlLadderHitCollection} \ref{sec:StMcPxlLadderHitCollection},
    \name{StMcPxlSensorHitCollection} \ref{sec:StMcPxlSensorHitCollection},
    and \name{StMcPxlHit} \ref{sec:StMcPxlHit}.
    \index{StMcPxlSectorHitCollection}
    \index{StMcPxlLadderHitCollection}
    \index{StMcPxlSensorHitCollection}
    \index{StMcPxlHit}

\item[Public\\ Constructors]
    \verb+StMcPxlHitCollection();+\\
    Create an instance of \name{StMcPxlHitCollection} and
    initialise the internal sector sub-collections.

\item[Public Member\\ Functions]

    \verb+bool addHit(StMcPxlHit*);+\\
    Add a hit to the appropriate sector sub-collection.

    \verb+unsigned int numberOfHits() const;+\\
    Count the total number of hits across all sectors.

    \verb+unsigned int numberOfSectors() const;+\\
    Return the number of sector sub-collections (10).

    \verb+StMcPxlSectorHitCollection*       sector(unsigned int);+\\
    \verb+const StMcPxlSectorHitCollection* sector(unsigned int) const;+\\
    Return the mutable or const sector sub-collection at 0-based index
    \name{i}.

\item[Examples]
{\footnotesize
\begin{verbatim}
//
// Count total PXL hits.
//
StMcPxlHitCollection* coll = mcEvt->pxlHitCollection();
cout << "Total PXL hits: " << coll->numberOfHits() << endl;
\end{verbatim}
}%footnotesize
\end{Entry}
\clearpage

%%%%%%%%%%%%%%%%%%%%%%%%%%%%%%%%%%%%%%%%%%%%%%%%%%%%%%%%%%%%%%%%%%%%
%
%    Reference: StMcPxlLadderHitCollection
%
%%%%%%%%%%%%%%%%%%%%%%%%%%%%%%%%%%%%%%%%%%%%%%%%%%%%%%%%%%%%%%%%%%%%
\subsection{StMcPxlLadderHitCollection}
\index{PXL ladder hit collection} \index{StMcPxlLadderHitCollection|textbf}
\label{sec:StMcPxlLadderHitCollection}
\begin{Entry}
\item[Summary]
    \name{StMcPxlLadderHitCollection} is a container for PXL hits in one
    ladder, organised by sensor.

\item[Synopsis]
    \verb+#include "StMcPxlLadderHitCollection.hh"+\\
    \verb+class StMcPxlLadderHitCollection;+\\

\item[Description]
    \name{StMcPxlLadderHitCollection} is the third level in the
    sector--ladder--sensor hit hierarchy for the PXL detector.  Each
    ladder holds an array of 10 \name{StMcPxlSensorHitCollection} objects,
    one per sensor.

\item[Persistence]
    None

\item[Related Classes]
    \name{StMcPxlLadderHitCollection} contains an array of
    {\tt StMcPxlSensorHitCollection} objects (size~10).  See
    \name{StMcPxlSectorHitCollection} \ref{sec:StMcPxlSectorHitCollection},
    \name{StMcPxlSensorHitCollection} \ref{sec:StMcPxlSensorHitCollection},
    and \name{StMcPxlHit} \ref{sec:StMcPxlHit}.
    \index{StMcPxlSectorHitCollection}
    \index{StMcPxlSensorHitCollection}
    \index{StMcPxlHit}

\item[Public\\ Constructors]
    \verb+StMcPxlLadderHitCollection();+\\
    Create an instance of \name{StMcPxlLadderHitCollection} and
    initialise the internal sensor sub-collections.

\item[Public Member\\ Functions]

    \verb+unsigned int numberOfHits() const;+\\
    Count the total number of hits across all sensors in this ladder.

    \verb+unsigned int numberOfSensors() const;+\\
    Return the number of sensor sub-collections (10).

    \verb+StMcPxlSensorHitCollection*       sensor(unsigned int);+\\
    \verb+const StMcPxlSensorHitCollection* sensor(unsigned int) const;+\\
    Return the mutable or const sensor sub-collection at 0-based index
    \name{i}.

\item[Examples]
{\footnotesize
\begin{verbatim}
//
// Look in StMcPxlSensorHitCollection for an example.
// These classes normally go hand in hand.
\end{verbatim}
}%footnotesize
\end{Entry}
\clearpage

%%%%%%%%%%%%%%%%%%%%%%%%%%%%%%%%%%%%%%%%%%%%%%%%%%%%%%%%%%%%%%%%%%%%
%
%    Reference: StMcPxlSectorHitCollection
%
%%%%%%%%%%%%%%%%%%%%%%%%%%%%%%%%%%%%%%%%%%%%%%%%%%%%%%%%%%%%%%%%%%%%
\subsection{StMcPxlSectorHitCollection}
\index{PXL sector hit collection} \index{StMcPxlSectorHitCollection|textbf}
\label{sec:StMcPxlSectorHitCollection}
\begin{Entry}
\item[Summary]
    \name{StMcPxlSectorHitCollection} is a container for PXL hits in one
    sector, organised by ladder.

\item[Synopsis]
    \verb+#include "StMcPxlSectorHitCollection.hh"+\\
    \verb+class StMcPxlSectorHitCollection;+\\

\item[Description]
    \name{StMcPxlSectorHitCollection} is the second level in the
    sector--ladder--sensor hit hierarchy for the PXL detector.  Each
    sector holds an array of 4 \name{StMcPxlLadderHitCollection} objects,
    one per ladder.

\item[Persistence]
    None

\item[Related Classes]
    \name{StMcPxlSectorHitCollection} contains an array of
    {\tt StMcPxlLadderHitCollection} objects (size~4).  See
    \name{StMcPxlHitCollection} \ref{sec:StMcPxlHitCollection},
    \name{StMcPxlLadderHitCollection} \ref{sec:StMcPxlLadderHitCollection},
    and \name{StMcPxlHit} \ref{sec:StMcPxlHit}.
    \index{StMcPxlHitCollection}
    \index{StMcPxlLadderHitCollection}
    \index{StMcPxlHit}

\item[Public\\ Constructors]
    \verb+StMcPxlSectorHitCollection();+\\
    Create an instance of \name{StMcPxlSectorHitCollection} and
    initialise the internal ladder sub-collections.

\item[Public Member\\ Functions]

    \verb+unsigned int numberOfHits() const;+\\
    Count the total number of hits across all ladders in this sector.

    \verb+unsigned int numberOfLadders() const;+\\
    Return the number of ladder sub-collections (4).

    \verb+StMcPxlLadderHitCollection*       ladder(unsigned int);+\\
    \verb+const StMcPxlLadderHitCollection* ladder(unsigned int) const;+\\
    Return the mutable or const ladder sub-collection at 0-based index
    \name{i}.

\item[Examples]
{\footnotesize
\begin{verbatim}
//
// Traverse the full PXL hierarchy to print all hit positions.
//
StMcPxlHitCollection* pxl = mcEvt->pxlHitCollection();
for (unsigned int s = 0; s < pxl->numberOfSectors(); ++s) {
  for (unsigned int ld = 0;
       ld < pxl->sector(s)->numberOfLadders(); ++ld) {
    for (unsigned int sn = 0;
         sn < pxl->sector(s)->ladder(ld)->numberOfSensors(); ++sn) {
      StSPtrVecMcPxlHit& hits =
          pxl->sector(s)->ladder(ld)->sensor(sn)->hits();
      for (StMcPxlHitIterator it = hits.begin();
           it != hits.end(); ++it)
          cout << (*it)->position() << endl;
    }
  }
}
\end{verbatim}
}%footnotesize
\end{Entry}
\clearpage

%%%%%%%%%%%%%%%%%%%%%%%%%%%%%%%%%%%%%%%%%%%%%%%%%%%%%%%%%%%%%%%%%%%%
%
%    Reference: StMcPxlSensorHitCollection
%
%%%%%%%%%%%%%%%%%%%%%%%%%%%%%%%%%%%%%%%%%%%%%%%%%%%%%%%%%%%%%%%%%%%%
\subsection{StMcPxlSensorHitCollection}
\index{PXL sensor hit collection} \index{StMcPxlSensorHitCollection|textbf}
\label{sec:StMcPxlSensorHitCollection}
\begin{Entry}
\item[Summary]
    \name{StMcPxlSensorHitCollection} is the leaf node in the PXL hit
    hierarchy; it directly stores the PXL hits for one sensor.

\item[Synopsis]
    \verb+#include "StMcPxlSensorHitCollection.hh"+\\
    \verb+class StMcPxlSensorHitCollection;+\\

\item[Description]
    \name{StMcPxlSensorHitCollection} stores the actual
    \name{StMcPxlHit} pointers in a \name{StSPtrVecMcPxlHit} container
    for a single PXL sensor.

\item[Persistence]
    None

\item[Related Classes]
    \name{StMcPxlSensorHitCollection} contains a
    {\tt StSPtrVecMcPxlHit}, which owns the hit objects.  See the
    section on containers \ref{sec:containers}.
    The parent collection is
    \name{StMcPxlLadderHitCollection} \ref{sec:StMcPxlLadderHitCollection}.
    \index{StMcPxlLadderHitCollection}
    \index{StMcPxlHit}
    \index{StSPtrVecMcPxlHit}

\item[Public\\ Constructors]
    \verb+StMcPxlSensorHitCollection();+\\
    Create an instance of \name{StMcPxlSensorHitCollection} and
    initialise the internal hit container.

\item[Public Member\\ Functions]

    \verb+StSPtrVecMcPxlHit&       hits();+\\
    \verb+const StSPtrVecMcPxlHit& hits() const;+\\
    Return a reference to the container of PXL hits for this sensor.

\item[Examples]
{\footnotesize
\begin{verbatim}
//
// Look in StMcPxlSectorHitCollection for a full traversal example.
\end{verbatim}
}%footnotesize
\end{Entry}
\clearpage

%%%%%%%%%%%%%%%%%%%%%%%%%%%%%%%%%%%%%%%%%%%%%%%%%%%%%%%%%%%%%%%%%%%%
%
%    Reference: StMcSsdHit
%
%%%%%%%%%%%%%%%%%%%%%%%%%%%%%%%%%%%%%%%%%%%%%%%%%%%%%%%%%%%%%%%%%%%%
\subsection{StMcSsdHit}
\index{SSD hit} \index{StMcSsdHit|textbf}
\label{sec:StMcSsdHit}
\begin{Entry}
\item[Summary]
    \name{StMcSsdHit} represents a monte carlo SSD (Silicon Strip Detector) hit.

\item[Synopsis]
    \verb+#include "StMcSsdHit.hh"+\\
    \verb+class StMcSsdHit;+\\

\item[Description]
    \name{StMcSsdHit} inherits most functionality from \name{StMcHit}.
    For a complete description of the inherited member functions see
    \name{StMcHit} (sec.~\ref{sec:StMcHit}).

    In the \StMcEvent\ data model each track keeps references to the
    associated hits and vice versa.  All hits have a pointer to their
    parent track.  This is especially important for the making of the
    associations in \StAssociationMaker.

\item[Persistence]
    None

\item[Related Classes]
    \name{StMcSsdHit} is derived directly from \name{StMcHit}.
    The way to retrieve hits for a given ladder and wafer starting
    from the top-level \name{StMcEvent} pointer is:\\
    \verb+mcEvt->ssdHitCollection()->ladder(ld)->wafer(w)->hits()+\\
    The hits are kept in an instance of \name{StSPtrVecMcSsdHit}
    where they are stored by pointer (see \ref{sec:containers}).
    Each instance of \name{StMcTrack} holds a list of SSD hits
    which belong to that track.
    \index{StMcHit}
    \index{StMcTrack}
    \index{StMcSsdHitCollection}

\item[Public\\ Constructors]
    \verb+StMcSsdHit();+\\
    Default constructor.

    \verb+StMcSsdHit(const StThreeVectorF& x, const StThreeVectorF& p,+\\
    \verb+           Float_t de, Float_t ds, Float_t tof,+\\
    \verb+           Long_t k, Long_t volId, StMcTrack* parent);+\\
    Create an instance of \name{StMcSsdHit} with position \name{x},
    local momentum \name{p}, energy deposition \name{de}, path length
    \name{ds}, time of flight \name{tof}, primary key \name{k},
    volume id \name{volId} and parent track \name{parent}.

    \verb+StMcSsdHit(g2t_ssd_hit_st* pt);+\\
    Construct from a GEANT \texttt{g2t\_ssd\_hit} table row.

\item[Public Member\\ Functions]

    \verb+ULong_t ladder() const;+\\
    Returns the ladder number decoded from the volume id.
    \index{ladder, SSD}

    \verb+ULong_t wafer() const;+\\
    Returns the wafer number decoded from the volume id.
    \index{wafer, SSD}

\item[Examples]
{\footnotesize
\begin{verbatim}
//
//  Print SSD hits belonging to a track.
//
void printSsdHits(StMcTrack* track)
{
    StPtrVecMcSsdHit& hits = track->ssdHits();
    for (StMcSsdHitIterator it = hits.begin();
         it != hits.end(); ++it)
        cout << *(*it) << endl;
}
\end{verbatim}
}%footnotesize
\end{Entry}
\clearpage

%%%%%%%%%%%%%%%%%%%%%%%%%%%%%%%%%%%%%%%%%%%%%%%%%%%%%%%%%%%%%%%%%%%%
%
%    Reference: StMcSsdHitCollection
%
%%%%%%%%%%%%%%%%%%%%%%%%%%%%%%%%%%%%%%%%%%%%%%%%%%%%%%%%%%%%%%%%%%%%
\subsection{StMcSsdHitCollection}
\index{SSD hit collection} \index{StMcSsdHitCollection|textbf}
\label{sec:StMcSsdHitCollection}
\begin{Entry}
\item[Summary]
    \name{StMcSsdHitCollection} is the top-level container for SSD monte
    carlo hits, organised by ladder.

\item[Synopsis]
    \verb+#include "StMcSsdHitCollection.hh"+\\
    \verb+class StMcSsdHitCollection;+\\

\item[Description]
    \name{StMcSsdHitCollection} is the first step down the hierarchy of hits
    for the SSD.  It provides methods to return a particular ladder
    sub-collection and to count the total number of hits across all SSD
    ladders.  The collection holds an array of 20
    \name{StMcSsdLadderHitCollection} objects, one per ladder.

\item[Persistence]
    None

\item[Related Classes]
    \name{StMcSsdHitCollection} contains an array of
    {\tt StMcSsdLadderHitCollection} objects (size~20).  See
    \name{StMcSsdLadderHitCollection} \ref{sec:StMcSsdLadderHitCollection},
    \name{StMcSsdWaferHitCollection} \ref{sec:StMcSsdWaferHitCollection},
    and \name{StMcSsdHit} \ref{sec:StMcSsdHit}.
    \index{StMcSsdLadderHitCollection}
    \index{StMcSsdWaferHitCollection}
    \index{StMcSsdHit}

\item[Public\\ Constructors]
    \verb+StMcSsdHitCollection();+\\
    Create an instance of \name{StMcSsdHitCollection} and
    initialise the internal ladder sub-collections.

\item[Public Member\\ Functions]

    \verb+bool addHit(StMcSsdHit*);+\\
    Add a hit to the appropriate ladder sub-collection.

    \verb+unsigned long numberOfHits() const;+\\
    Count the total number of hits across all ladders.

    \verb+unsigned int  numberOfLadders() const;+\\
    Return the number of ladder sub-collections (20).

    \verb+StMcSsdLadderHitCollection*       ladder(unsigned int);+\\
    \verb+const StMcSsdLadderHitCollection* ladder(unsigned int) const;+\\
    Return the mutable or const ladder sub-collection at 0-based index
    \name{i}.

\item[Examples]
{\footnotesize
\begin{verbatim}
//
// Count total SSD hits.
//
StMcSsdHitCollection* coll = mcEvt->ssdHitCollection();
cout << "Total SSD hits: " << coll->numberOfHits() << endl;
\end{verbatim}
}%footnotesize
\end{Entry}
\clearpage

%%%%%%%%%%%%%%%%%%%%%%%%%%%%%%%%%%%%%%%%%%%%%%%%%%%%%%%%%%%%%%%%%%%%
%
%    Reference: StMcSsdLadderHitCollection
%
%%%%%%%%%%%%%%%%%%%%%%%%%%%%%%%%%%%%%%%%%%%%%%%%%%%%%%%%%%%%%%%%%%%%
\subsection{StMcSsdLadderHitCollection}
\index{SSD ladder hit collection} \index{StMcSsdLadderHitCollection|textbf}
\label{sec:StMcSsdLadderHitCollection}
\begin{Entry}
\item[Summary]
    \name{StMcSsdLadderHitCollection} is a container for SSD hits in one
    ladder, organised by wafer.

\item[Synopsis]
    \verb+#include "StMcSsdLadderHitCollection.hh"+\\
    \verb+class StMcSsdLadderHitCollection;+\\

\item[Description]
    \name{StMcSsdLadderHitCollection} is the second level in the
    ladder--wafer hit hierarchy for the SSD.  Each ladder holds an array of
    up to 16 \name{StMcSsdWaferHitCollection} objects.

\item[Persistence]
    None

\item[Related Classes]
    \name{StMcSsdLadderHitCollection} contains an array of
    {\tt StMcSsdWaferHitCollection} objects (size~16).  See
    \name{StMcSsdHitCollection} \ref{sec:StMcSsdHitCollection},
    \name{StMcSsdWaferHitCollection} \ref{sec:StMcSsdWaferHitCollection},
    and \name{StMcSsdHit} \ref{sec:StMcSsdHit}.
    \index{StMcSsdHitCollection}
    \index{StMcSsdWaferHitCollection}
    \index{StMcSsdHit}

\item[Public\\ Constructors]
    \verb+StMcSsdLadderHitCollection();+\\
    Create an instance of \name{StMcSsdLadderHitCollection} and
    initialise the internal wafer sub-collections.

\item[Public Member\\ Functions]

    \verb+unsigned long numberOfHits() const;+\\
    Count the total number of hits across all wafers in this ladder.

    \verb+unsigned int  numberOfWafers() const;+\\
    Return the number of wafer sub-collections (16).

    \verb+StMcSsdWaferHitCollection*       wafer(unsigned int);+\\
    \verb+const StMcSsdWaferHitCollection* wafer(unsigned int) const;+\\
    Return the mutable or const wafer sub-collection at 0-based index
    \name{i}.

\item[Examples]
{\footnotesize
\begin{verbatim}
//
// Look in StMcSsdWaferHitCollection for an example.
// These classes normally go hand in hand.
\end{verbatim}
}%footnotesize
\end{Entry}
\clearpage

%%%%%%%%%%%%%%%%%%%%%%%%%%%%%%%%%%%%%%%%%%%%%%%%%%%%%%%%%%%%%%%%%%%%
%
%    Reference: StMcSsdWaferHitCollection
%
%%%%%%%%%%%%%%%%%%%%%%%%%%%%%%%%%%%%%%%%%%%%%%%%%%%%%%%%%%%%%%%%%%%%
\subsection{StMcSsdWaferHitCollection}
\index{SSD wafer hit collection} \index{StMcSsdWaferHitCollection|textbf}
\label{sec:StMcSsdWaferHitCollection}
\begin{Entry}
\item[Summary]
    \name{StMcSsdWaferHitCollection} is the leaf node in the SSD hit
    hierarchy; it directly stores the SSD hits for one wafer.

\item[Synopsis]
    \verb+#include "StMcSsdWaferHitCollection.hh"+\\
    \verb+class StMcSsdWaferHitCollection;+\\

\item[Description]
    \name{StMcSsdWaferHitCollection} stores the actual
    \name{StMcSsdHit} pointers in a \name{StSPtrVecMcSsdHit} container
    for a single SSD wafer.

\item[Persistence]
    None

\item[Related Classes]
    \name{StMcSsdWaferHitCollection} contains a
    {\tt StSPtrVecMcSsdHit}, which owns the hit objects.  See the
    section on containers \ref{sec:containers}.
    The parent collection is
    \name{StMcSsdLadderHitCollection} \ref{sec:StMcSsdLadderHitCollection}.
    \index{StMcSsdLadderHitCollection}
    \index{StMcSsdHit}
    \index{StSPtrVecMcSsdHit}

\item[Public\\ Constructors]
    \verb+StMcSsdWaferHitCollection();+\\
    Create an instance of \name{StMcSsdWaferHitCollection} and
    initialise the internal hit container.

\item[Public Member\\ Functions]

    \verb+StSPtrVecMcSsdHit&       hits();+\\
    \verb+const StSPtrVecMcSsdHit& hits() const;+\\
    Return a reference to the container of SSD hits for this wafer.

\item[Examples]
{\footnotesize
\begin{verbatim}
//
// Print number of SSD hits and their positions.
//
StMcSsdHitCollection* ssdColl = mcEvt->ssdHitCollection();
cout << "Total SSD hits: " << ssdColl->numberOfHits() << endl;
for (unsigned int ld = 0; ld < ssdColl->numberOfLadders(); ++ld) {
  for (unsigned int w = 0;
       w < ssdColl->ladder(ld)->numberOfWafers(); ++w) {
    StSPtrVecMcSsdHit& hits =
        ssdColl->ladder(ld)->wafer(w)->hits();
    for (StMcSsdHitIterator it = hits.begin();
         it != hits.end(); ++it)
        cout << (*it)->position() << endl;
  }
}
\end{verbatim}
}%footnotesize
\end{Entry}
\clearpage

%%%%%%%%%%%%%%%%%%%%%%%%%%%%%%%%%%%%%%%%%%%%%%%%%%%%%%%%%%%%%%%%%%%%
%
%    Reference: StMcSvtBarrelHitCollection
%
%%%%%%%%%%%%%%%%%%%%%%%%%%%%%%%%%%%%%%%%%%%%%%%%%%%%%%%%%%%%%%%%%%%%
\subsection{StMcSvtBarrelHitCollection}
\index{SVT barrel hit collection} \index{StMcSvtBarrelHitCollection|textbf}
\label{sec:StMcSvtBarrelHitCollection}
\textbf{[Renamed from StMcSvtLayerHitCollection]}
This class replaced \textbf{StMcSvtLayerHitCollection} when the SVT hit
storage scheme was reorganised from layer--ladder--wafer to
barrel--ladder--wafer order (commit \texttt{f98478bb68}).
The stub entry for the old name is at \S\ref{sec:StMcSvtLayerHitCollection}.
\begin{Entry}
\item[Summary]
    \name{StMcSvtBarrelHitCollection} is a container for SVT hits in one
    barrel, organised by ladder.  It is the second level in the
    barrel--ladder--wafer hit hierarchy.

\item[Synopsis]
    \verb+#include "StMcSvtBarrelHitCollection.hh"+\\
    \verb+class StMcSvtBarrelHitCollection;+\\

\item[Description]
    \name{StMcSvtBarrelHitCollection} is the second step down the
    hierarchy of hits for the SVT.  Each barrel holds an array of up to 20
    \name{StMcSvtLadderHitCollection} objects.  The capacity is 20 to
    accommodate the SSD barrel, which has more ladders than each SVT barrel.

    This class replaces \name{StMcSvtLayerHitCollection}
    (see \S\ref{sec:StMcSvtLayerHitCollection}), reflecting the change in
    the SVT storage scheme from layer--ladder--wafer to
    barrel--ladder--wafer introduced in commit \texttt{f98478bb68}.

\item[Persistence]
    None

\item[Related Classes]
    \name{StMcSvtBarrelHitCollection} contains an array of
    {\tt StMcSvtLadderHitCollection} objects (size~20).  See
    \name{StMcSvtHitCollection} \ref{sec:StMcSvtHitCollection},
    \name{StMcSvtLadderHitCollection} \ref{sec:StMcSvtLadderHitCollection},
    and \name{StMcSvtWaferHitCollection} \ref{sec:StMcSvtWaferHitCollection}.
    \index{StMcSvtHitCollection}
    \index{StMcSvtLadderHitCollection}
    \index{StMcSvtWaferHitCollection}
    \index{StMcSvtHit}

\item[Public\\ Constructors]
    \verb+StMcSvtBarrelHitCollection();+\\
    Create an instance of \name{StMcSvtBarrelHitCollection} and
    initialise the internal ladder sub-collections.

\item[Public Member\\ Functions]

    \verb+unsigned long numberOfHits() const;+\\
    Count the total number of hits across all ladders in this barrel.

    \verb+unsigned int  numberOfLadders() const;+\\
    Return the number of ladder sub-collections (up to 20).

    \verb+StMcSvtLadderHitCollection*       ladder(unsigned int);+\\
    \verb+const StMcSvtLadderHitCollection* ladder(unsigned int) const;+\\
    Return the mutable or const ladder sub-collection at 0-based index
    \name{i}.

    \verb+void setBarrelNumber(int);+\\
    Set the barrel index and propagate it to all ladder sub-collections.

\item[Examples]
{\footnotesize
\begin{verbatim}
//
// Look in StMcSvtWaferHitCollection for a full traversal example.
// These classes normally go hand in hand.
\end{verbatim}
}%footnotesize
\end{Entry}
\clearpage

\printindex

\end{document}
\bye
% The text following after this line is not included into the text

%
%  Template for reference section
%

%%%%%%%%%%%%%%%%%%%%%%%%%%%%%%%%%%%%%%%%%%%%%%%%%%%%%%%%%%%%%%%%%%%%
%
%    Reference: className
%
%%%%%%%%%%%%%%%%%%%%%%%%%%%%%%%%%%%%%%%%%%%%%%%%%%%%%%%%%%%%%%%%%%%%
\subsection{className}
\index{className|textbf}
\label{sec:className}
\begin{Entry}
\item[Summary]

\item[Synopsis]
    \verb+#include "className.hh"+\\
    \verb+class className;+\\

\item[Description]

\item[Persistence]
    None

\item[Related Classes]

\item[Public\\ Constructors]

\item[Public Member\\ Functions]

\item[Public Member\\ Operators]

\item[Public Functions]

\item[Public Operators]

\item[Examples]
{\footnotesize
\begin{verbatim}
//
//  What the example does
//

\end{verbatim}
}%footnotesize

\end{Entry}
